\section{Huyệt}

\section{LU 7 -- Liệt khuyết (Lieque)}

\subsection{Vị trí huyệt}

Huyệt Liệt khuyết nằm ở mặt quay của cẳng tay, cách nếp gấp ngang cổ tay khoảng \(1{,}5\) thốn về phía gần khuỷu tay, trong khe giữa gân cơ cánh tay quay và gân cơ dạng dài ngón cái.

Có thể xác định huyệt bằng cách đan chéo ngón trỏ và ngón cái của hai bàn tay, đặt ngón trỏ của một bàn tay lên mỏm trâm quay của bàn tay còn lại. Huyệt nằm tại chỗ lõm ngay dưới đầu ngón trỏ.

\begin{figure}[htbp]
    \centering
    \includegraphics[width=0.65\textwidth]{LU_7.jpg}
    \caption{Huyệt Liệt khuyết (LU 7).}
    \label{fig:lu7}
\end{figure}

\subsection{Chỉ định}

Các chỉ định được liệt kê trong tài liệu nguồn bao gồm:

\begin{itemize}
    \item Ho, hen suyễn và đau họng.
    \item Đau đầu, đau nửa đầu và cứng cổ.
    \item Liệt mặt và đau răng.
    \item Đau và yếu vùng cổ tay.
\end{itemize}

\subsection{Cơ chế tác dụng theo y học cổ truyền}

Theo lý luận y học cổ truyền, huyệt Liệt khuyết có các tác dụng sau:

\begin{itemize}
    \item Giải biểu, khu phong.
    \item Tuyên thông Phế khí, thúc đẩy chức năng túc giáng của Phế.
    \item Khai thông và điều hòa mạch Nhâm.
    \item Điều hòa thủy đạo.
    \item Có lợi cho vùng đầu và cổ.
\end{itemize}

\subsection{Ứng dụng trong phương pháp Balance}

Trong phương pháp Balance, huyệt Liệt khuyết được liên hệ với các kinh mạch và hệ thống sau:

\begin{itemize}
    \item \textbf{Kinh thủ Thái âm -- Phế:} kinh mạch chủ quản của huyệt.
    \item \textbf{Hệ thống 1:} kinh túc Thái âm -- Tỳ.
    \item \textbf{Hệ thống 2:} kinh túc Thái dương -- Bàng quang.
    \item \textbf{Hệ thống 3:} kinh thủ Dương minh -- Đại trường.
    \item \textbf{Hệ thống 4:} kinh túc Thái dương -- Bàng quang.
    \item \textbf{Hệ thống 5:} kinh túc Quyết âm -- Can.
\end{itemize}

\subsection{Kỹ thuật châm cứu}

Theo tài liệu nguồn, kỹ thuật châm được mô tả như sau:

\begin{itemize}
    \item \textbf{Hướng châm:} châm chếch hướng lên phía khuỷu tay.
    \item \textbf{Độ sâu:} khoảng \(0{,}3\)--\(0{,}5\) thốn.
\end{itemize}

Hướng châm và độ sâu thực tế cần được điều chỉnh theo đặc điểm giải phẫu, thể trạng người bệnh và đánh giá của người hành nghề có chuyên môn.

\subsection{Đặc tính huyệt}

Theo phân loại trong tài liệu nguồn, Liệt khuyết có các đặc tính sau:

\begin{itemize}
    \item \textbf{Lạc huyệt:} kinh Phế.
    \item \textbf{Huyệt Xuất:} kinh Phế.
    \item \textbf{Huyệt phối:} mạch Âm Kiều (Yinqiao).
    \item \textbf{Huyệt chủ:} mạch Nhâm.
    \item \textbf{Huyệt chỉ huy:} vùng đầu và gáy.
\end{itemize}

\subsection{Tài liệu tham khảo}

Master Tung's Acupuncture. (n.d.). \textit{Lieque (LU 7)}.
eLotus CORE.
\url{https://www.mastertungacupuncture.org/acupuncture/traditional/points/lu7}
% 
\clearpage
\subsection{Lung Meridian (LU): Kinh thủ thái âm Phế}
\subsubsection{Zhongfu (LU 1): Trung Phủ}
\includegraphics[width=\linewidth]{"LU_1.jpg"}

\subsubsection{Yunmen (LU 2): Vân Môn}
\includegraphics[width=\linewidth]{"LU_2.jpg"}

\subsubsection{Tianfu (LU 3): Thiên Phủ}
\includegraphics[width=\linewidth]{"LU_3.jpg"}

\subsubsection{Xiabai (LU 4): Hiệp Bạch}
\includegraphics[width=\linewidth]{"LU_4.jpg"}

\subsubsection{Chize (LU 5): Xích Trạch}
\includegraphics[width=\linewidth]{"LU_5.jpg"}

\subsubsection{Kongzui (LU 6): Khổng Tối}
\includegraphics[width=\linewidth]{"LU_6.jpg"}

\subsubsection{Lieque (LU 7): Liệt Khuyết}
\includegraphics[width=\linewidth]{"LU_7.jpg"}

\subsubsection{Jingqu (LU 8): Kinh Cừ}
\includegraphics[width=\linewidth]{"LU_8.jpg"}

\subsubsection{Taiyuan (LU 9): Thái Uyên}
\includegraphics[width=\linewidth]{"LU_9.jpg"}

\subsubsection{Yuji (LU 10): Ngư Tế}
\includegraphics[width=\linewidth]{"LU_10.jpg"}

\subsubsection{Shaoshang (LU 11): Thiếu Thương}
\includegraphics[width=\linewidth]{"LU_11.jpg"}

% \clearpage
\subsection{Large Intestine Meridian (LI): Kinh thủ dương minh Đại trường}

\subsubsection{Shangyang (LI 1): Thương Dương}
\includegraphics[width=\linewidth]{"LI_1.jpg"}

\subsubsection{Erjian (LI 2): Nhị Gian}
\includegraphics[width=\linewidth]{"LI_2.jpg"}

\subsubsection{Sanjian (LI 3): Tam Gian}
\includegraphics[width=\linewidth]{"LI_3.jpg"}

\subsubsection{Hegu (LI 4): Hợp Cốc}
\includegraphics[width=\linewidth]{"LI_4.jpg"}

\subsubsection{Yangxi (LI 5): Dương Khê}
\includegraphics[width=\linewidth]{"LI_5.jpg"}

\subsubsection{Pianli (LI 6): Thiên Lịch}
\includegraphics[width=\linewidth]{"LI_6.jpg"}

\subsubsection{Wenliu (LI 7): Ôn Lưu}
\includegraphics[width=\linewidth]{"LI_7.jpg"}

\subsubsection{Xialian (LI 8): Hạ Liêm}
\includegraphics[width=\linewidth]{"LI_8.jpg"}

\subsubsection{Shanglian (LI 9): Thượng Liêm}
\includegraphics[width=\linewidth]{"LI_9.jpg"}

\subsubsection{Shousanli (LI 10): Thủ Tam Lý}
\includegraphics[width=\linewidth]{"LI_10.jpg"}

\subsubsection{Quchi (LI 11): Khúc Trì}
\includegraphics[width=\linewidth]{"LI_11.jpg"}

\subsubsection{Zhouliao (LI 12): Trửu Liêu}
\includegraphics[width=\linewidth]{"LI_12a_0.jpg"}

\subsubsection{Shouwuli (LI 13): Thủ Ngũ Lý}
\includegraphics[width=\linewidth]{"LI_13.jpg"}

\subsubsection{Binao (LI 14): Tý Nhu}
\includegraphics[width=\linewidth]{"LI_14.jpg"}

\subsubsection{Jianyu (LI 15): Kiên Ngung}
\includegraphics[width=\linewidth]{"LI_15b.jpg"}

\subsubsection{Jugu (LI 16): Cự Cốt}
\includegraphics[width=\linewidth]{"LI_16b.jpg"}

\subsubsection{Tianding (LI 17): Thiên Đỉnh}
\includegraphics[width=\linewidth]{"LI 17.jpg"}

\subsubsection{Futu (LI 18): Phủ Đột}
\includegraphics[width=\linewidth]{"LI 18.jpg"}

\subsubsection{Kouheliao (LI 19): Hòa Liêu}
\includegraphics[width=\linewidth]{"LI 19.jpg"}

\subsubsection{Yingxiang (LI 20): Nghinh Hương}
\includegraphics[width=\linewidth]{"LI 20.jpg"}


% \clearpage
\subsection{Stomach Meridian (ST): Kinh túc dương minh Vị}

\subsubsection{Chengqi (ST 1): Thừa Khấp}
\includegraphics[width=\linewidth]{"ST 1.jpg"}
Huyệt này nằm ngay phía dưới đồng tử (khi mắt nhìn thẳng về phía trước),
ở khoảng giữa nhãn cầu và bờ dưới hốc mắt.

\subsubsection{Sibai (ST 2): Tứ Bạch}
\includegraphics[width=\linewidth]{"ST 2.jpg"}

\subsubsection{Juliao (ST 3): Cự Liêu}
\includegraphics[width=\linewidth]{"ST 3.jpg"}

\subsubsection{Dicang (ST 4): Địa Thương}
\includegraphics[width=\linewidth]{"ST 4.jpg"}

\subsubsection{Daying (ST 5): Đại Nghinh}
\includegraphics[width=\linewidth]{"ST 5.jpg"}

\subsubsection{Jiache (ST 6): Giáp Xa}
\includegraphics[width=\linewidth]{"ST 6_1_0.jpg"}
\includegraphics[width=\linewidth]{"ST 6_2.jpg"}

\subsubsection{Xiaguan (ST 7): Hạ Quan}
\includegraphics[width=\linewidth]{"ST 7.jpg"}

\subsubsection{Touwei (ST 8): Đầu Duy}
\includegraphics[width=\linewidth]{"ST 8.jpg"}
\includegraphics[width=\linewidth]{"ST 8_2.jpg"}

\subsubsection{Renying (ST 9): Nhân Nghinh}
\includegraphics[width=\linewidth]{"ST 9.jpg"}

\subsubsection{Shuitu (ST 10): Thủy Đột}
\includegraphics[width=\linewidth]{"ST 10.jpg"}

\subsubsection{Qishe (ST 11): Khí Xá}
\includegraphics[width=\linewidth]{"ST 11.jpg"}

\subsubsection{Quepen (ST 12): Khuyết Bồn}
\includegraphics[width=\linewidth]{"ST 12.jpg"}

\subsubsection{Qihu (ST 13): Khí Hộ}
\includegraphics[width=\linewidth]{"ST 13.jpg"}

\subsubsection{Kufang (ST 14): Khố Phòng}
\includegraphics[width=\linewidth]{"ST 14.jpg"}

\subsubsection{Wuyi (ST 15): Ốc É}
\includegraphics[width=\linewidth]{"ST 15.jpg"}

\subsubsection{Yingchuang (ST 16): Ưng Song}
\includegraphics[width=\linewidth]{"ST 16.jpg"}

\subsubsection{Ruzhong (ST 17): Nhũ Trung}
\includegraphics[width=\linewidth]{"ST 17.jpg"}

\subsubsection{Rugen (ST 18): Nhũ Căn}
\includegraphics[width=\linewidth]{"ST 18.jpg"}

\subsubsection{Burong (ST 19): Bất Dung}
\includegraphics[width=\linewidth]{"ST 19.jpg"}

\subsubsection{Chengman (ST 20): Thừa Mãn}
\includegraphics[width=\linewidth]{"ST 20.jpg"}

\subsubsection{Liangmen (ST 21): Lương Môn}
\includegraphics[width=\linewidth]{"ST 21.jpg"}

\subsubsection{Guanmen (ST 22): Quan Môn}
\includegraphics[width=\linewidth]{"ST 22.jpg"}

\subsubsection{Taiyi (ST 23): Thái Ất}
\includegraphics[width=\linewidth]{"ST 23.jpg"}

\subsubsection{Huaroumen (ST 24): Hoạt Nhục Môn}
\includegraphics[width=\linewidth]{"ST 24.jpg"}

\subsubsection{Tianshu (ST 25): Thiên Khu}
\includegraphics[width=\linewidth]{"ST 25.jpg"}

\subsubsection{Wailing (ST 26): Ngoại Lăng}
\includegraphics[width=\linewidth]{"ST 26.jpg"}

\subsubsection{Daju (ST 27): Đại Cự}
\includegraphics[width=\linewidth]{"ST 27.jpg"}

\subsubsection{Shuidao (ST 28): Thủy Đạo}
\includegraphics[width=\linewidth]{"ST 28.jpg"}

\subsubsection{Guilai (ST 29): Quy Lai}
\includegraphics[width=\linewidth]{"ST 29.jpg"}

\subsubsection{Qichong (ST 30): Khí Xung}
\includegraphics[width=\linewidth]{"ST 30.jpg"}

\subsubsection{Biguan (ST 31): Bễ Quan}
\includegraphics[width=\linewidth]{"ST_31.jpg"}

\subsubsection{Futu (ST 32): Phục Thỏ}
\includegraphics[width=\linewidth]{"ST_32.jpg"}

\subsubsection{Yinshi (ST 33): Âm Thị}
\includegraphics[width=\linewidth]{"ST_33.jpg"}

\subsubsection{Liangqiu (ST 34): Lương Khâu}
\includegraphics[width=\linewidth]{"ST_34.jpg"}

\subsubsection{Dubi (ST 35): Độc Tỵ}
\includegraphics[width=\linewidth]{"ST_35.jpg"}

\subsubsection{Zusanli (ST 36): Túc Tam Lý}
\includegraphics[width=\linewidth]{"ST_36.jpg"}

\subsubsection{Shangjuxu (ST 37): Thượng Cự Hư}
\includegraphics[width=\linewidth]{"ST_37.jpg"}

\subsubsection{Tiaokou (ST 38): Điều Khẩu}
\includegraphics[width=\linewidth]{"ST_38.jpg"}

\subsubsection{Xiajuxu (ST 39): Hạ Cự Hư}
\includegraphics[width=\linewidth]{"ST_39.jpg"}

\subsubsection{Fenglong (ST 40): Phong Long}
\includegraphics[width=\linewidth]{"ST_40.jpg"}

\subsubsection{Jiexi (ST 41): Giải Khê}
\includegraphics[width=\linewidth]{"ST 41.jpg"}

\subsubsection{Chongyang (ST 42): Xung Dương}
\includegraphics[width=\linewidth]{"ST 42.jpg"}

\subsubsection{Xiangu (ST 43): Hãm Cốc}
\includegraphics[width=\linewidth]{"ST_43.jpg"}

\subsubsection{Neiting (ST 44): Nội Đình}
\includegraphics[width=\linewidth]{"ST_44.jpg"}

\subsubsection{Lidui (ST 45): Lệ Đoài}
\includegraphics[width=\linewidth]{"ST_45.jpg"}

% \clearpage
\subsection{Spleen Meridian(SP): Kinh túc thái âm Tỳ}

\subsubsection{Yinbai (SP 1): Ân Bạch}
\includegraphics[width=\linewidth]{"SP_1.jpg"}

\subsubsection{Dadu (SP 2): Đại Đô}
\includegraphics[width=\linewidth]{"SP_2.jpg"}

\subsubsection{Taibai (SP 3): Thái Bạch}
\includegraphics[width=\linewidth]{"SP_3.jpg"}

\subsubsection{Gongsun (SP 4): Công Tôn}
\includegraphics[width=\linewidth]{"SP_4.jpg"}

\subsubsection{Shangqiu (SP 5): Thương Khâu}
\includegraphics[width=\linewidth]{"SP_5.jpg"}

\subsubsection{Sanyinjiao (SP 6): Tam Âm Giao}
\includegraphics[width=\linewidth]{"SP_6.jpg"}

\subsubsection{Lougu (SP 7): Lậu Cốc}
\includegraphics[width=\linewidth]{"SP_7.jpg"}

\subsubsection{Diji (SP 8): Địa Cơ}
\includegraphics[width=\linewidth]{"SP_8.jpg"}

\subsubsection{Yinlingquan (SP 9): Âm Lăng Tuyền}
\includegraphics[width=\linewidth]{"SP_9.jpg"}

\subsubsection{Xuehai (SP 10): Huyết Hải}
\includegraphics[width=\linewidth]{"SP_10a.jpg"}

\subsubsection{Jimen (SP 11): Cơ Môn}
\includegraphics[width=\linewidth]{"SP_11.jpg"}

\subsubsection{Chongmen (SP 12): Xung Môn}
\includegraphics[width=\linewidth]{"SP 12.jpg"}

\subsubsection{Fushe (SP 13): Phủ Xá}
\includegraphics[width=\linewidth]{"SP 13.jpg"}

\subsubsection{Fujie (SP 14): Phúc Kết}
\includegraphics[width=\linewidth]{"SP 14.jpg"}

\subsubsection{Daheng (SP 15): Đại Hoành}
\includegraphics[width=\linewidth]{"SP 15.jpg"}

\subsubsection{Fuai (SP 16): Phúc Ai}
\includegraphics[width=\linewidth]{"SP 16.jpg"}

\subsubsection{Shidou (SP 17): Thực Đậu}
\includegraphics[width=\linewidth]{"SP 17.jpg"}

\subsubsection{Tianxi (SP 18): Thiên Khê}
\includegraphics[width=\linewidth]{"SP 18.jpg"}

\subsubsection{Xiongxiang (SP 19): Hung Hương}
\includegraphics[width=\linewidth]{"SP 19.jpg"}

\subsubsection{Zhourong (SP 20): Chu Vinh}
\includegraphics[width=\linewidth]{"SP 20.jpg"}

\subsubsection{Dabao (SP 21): Đại Bao}
\includegraphics[width=\linewidth]{"SP21.jpg"}

% \clearpage
\subsection{Heart Meridian (HT): Kinh thủ thiếu âm Tâm}

\subsubsection{Jiquan (HT 1): Cực Tuyền}
\includegraphics[width=\linewidth]{"HT1.jpg"}
\includegraphics[width=\linewidth]{"HT 1.jpg"}

\subsubsection{Qingling (HT 2): Thanh Linh}
\includegraphics[width=\linewidth]{"HT_2.jpg"}

\subsubsection{Shaohai (HT 3): Thiếu Hải}
\includegraphics[width=\linewidth]{"HT_3.jpg"}

\subsubsection{Lingdao (HT 4): Linh Đạo}
\includegraphics[width=\linewidth]{"HT_4.jpg"}

\subsubsection{Tongli (HT 5): Thông Lý}
\includegraphics[width=\linewidth]{"HT_5.jpg"}

\subsubsection{Yinxi (HT 6): Âm Khích}
\includegraphics[width=\linewidth]{"HT_6.jpg"}

\subsubsection{Shenmen (HT 7): Thần Môn}
\includegraphics[width=\linewidth]{"HT_7a.jpg"}
\includegraphics[width=\linewidth]{"HT_7b.jpg"}

\subsubsection{Shaofu (HT 8): Thiếu Phủ}
\includegraphics[width=\linewidth]{"HT_8.jpg"}

\subsubsection{Shaochong (HT 9): Thiếu Xung}
\includegraphics[width=\linewidth]{"HT_9.jpg"}
% \clearpage
\subsection{Small Intestine Meridian (SI): Kinh thủ thái dương Tiểu trường}

\subsubsection{Shaoze (SI 1): Thiếu Trạch}
\includegraphics[width=\linewidth]{"SI_1b.jpg"}

\subsubsection{Qiangu (SI 2): Tiền Cốc}
\includegraphics[width=\linewidth]{"SI_2.jpg"}

\subsubsection{Houxi (SI 3): Hậu Khê}
\includegraphics[width=\linewidth]{"SI_3.jpg"}

\subsubsection{Wangu (SI 4): Uyển Cốt}
\includegraphics[width=\linewidth]{"SI_4.jpg"}

\subsubsection{Yanggu (SI 5): Dương Cốc}
\includegraphics[width=\linewidth]{"SI_5.jpg"}

\subsubsection{Yanglao (SI 6): Dưỡng Lão}
\includegraphics[width=\linewidth]{"SI_6.jpg"}

\subsubsection{Zhizheng (SI 7): Chi Chính}
\includegraphics[width=\linewidth]{"SI_7.jpg"}

\subsubsection{Xiaohai (SI 8): Tiểu Hải}
\includegraphics[width=\linewidth]{"SI_8.jpg"}

\subsubsection{Jianzhen (SI 9): Kiên Trinh}
\includegraphics[width=\linewidth]{"SI_9.jpg"}

\subsubsection{Naoshu (SI 10): Nhu Du}
\includegraphics[width=\linewidth]{"SI_10.jpg"}

\subsubsection{Tianzong (SI 11): Thiên Tông}
\includegraphics[width=\linewidth]{"SI_11.jpg"}

\subsubsection{Bingfeng (SI 12): Binh Phong}
\includegraphics[width=\linewidth]{"SI_12.jpg"}

\subsubsection{Quyuan (SI 13): Khúc Viên}
\includegraphics[width=\linewidth]{"SI_13.jpg"}

\subsubsection{Jianwaishu (SI 14): Kiên Ngoại Du}
\includegraphics[width=\linewidth]{"SI_14.jpg"}

\subsubsection{Jianzhongshu (SI 15): Kiên Trung Du}
\includegraphics[width=\linewidth]{"SI_15.jpg"}

\subsubsection{Tianchuang (SI 16): Thiên Song}
\includegraphics[width=\linewidth]{"SI 16.jpg"}

\subsubsection{Tianrong (SI 17): Thiên Dung}
\includegraphics[width=\linewidth]{"SI 17.jpg"}

\subsubsection{Quanliao (SI 18): Quyền Liêu}
\includegraphics[width=\linewidth]{"SI 18_0.jpg"}

\subsubsection{Tinggong (SI 19): Thính Cung}
\includegraphics[width=\linewidth]{"SI 19.jpg"}

% \clearpage
\subsection{Bladder Meridian (BL): Kinh thủ thái dương Bàng quang}

\subsubsection{Jingming (BL 1): Tình Minh}
\includegraphics[width=\linewidth]{"BL 1.jpg"}

\subsubsection{Zanzhu (BL 2): Toản Trúc}
\includegraphics[width=\linewidth]{"BL 2.jpg"}

\subsubsection{Meichong (BL 3): My Xung}
\includegraphics[width=\linewidth]{"B3_0.jpg"}

\subsubsection{Quchai (BL 4): Khúc Sai}
\includegraphics[width=\linewidth]{"BL4.jpg"}

\subsubsection{Wuchu (BL 5): Ngũ Xứ}
\includegraphics[width=\linewidth]{"BL5.jpg"}

\subsubsection{Chengguang (BL 6): Thừa Quang}
\includegraphics[width=\linewidth]{"BL6.jpg"}

\subsubsection{Tongtian (BL 7): Thông Thiên}
\includegraphics[width=\linewidth]{"BL7.jpg"}

\subsubsection{Luoque (BL 8): Lạc Khước}
\includegraphics[width=\linewidth]{"BL8.jpg"}

\subsubsection{Yuzhen (BL 9): Ngọc Chẩm}
\includegraphics[width=\linewidth]{"BL9.jpg"}

\subsubsection{Tianzhu (BL 10): Thiên Trụ}
\includegraphics[width=\linewidth]{"BL10.jpg"}

\subsubsection{Dazhu (BL 11): Đại Trữ}
\includegraphics[width=\linewidth]{"BL_11.jpg"}

\subsubsection{Fengmen (BL 12): Phong Môn}
\includegraphics[width=\linewidth]{"BL_12.jpg"}

\subsubsection{Feishu (BL 13): Phế Du}
\includegraphics[width=\linewidth]{"BL_13.jpg"}

\subsubsection{Jueyinshu (BL 14): Quyết Âm Du}
\includegraphics[width=\linewidth]{"BL_14.jpg"}

\subsubsection{Xinshu (BL 15): Tâm Du}
\includegraphics[width=\linewidth]{"BL_15.jpg"}

\subsubsection{Dushu (BL 16): Đốc Du}
\includegraphics[width=\linewidth]{"BL_16.jpg"}

\subsubsection{Geshu (BL 17): Cách Du}
\includegraphics[width=\linewidth]{"BL_17.jpg"}

\subsubsection{Ganshu (BL 18): Can Du}
\includegraphics[width=\linewidth]{"BL_18.jpg"}

\subsubsection{Danshu (BL 19): Đởm Du}
\includegraphics[width=\linewidth]{"BL_19.jpg"}

\subsubsection{Pishu (BL 20): Tỳ Du}
\includegraphics[width=\linewidth]{"BL_20.jpg"}

\subsubsection{Weishu (BL 21): Vị Du}
\includegraphics[width=\linewidth]{"BL_21.jpg"}

\subsubsection{Sanjiaoshu (BL 22): Tam Tiêu Du}
\includegraphics[width=\linewidth]{"BL_22.jpg"}

\subsubsection{Shenshu (BL 23): Thận Du}
\includegraphics[width=\linewidth]{"BL_23.jpg"}

\subsubsection{Qihaishu (BL 24): Khí Hải Du}
\includegraphics[width=\linewidth]{"BL_24.jpg"}

\subsubsection{Dachangshu (BL 25): Đại Trường Du}
\includegraphics[width=\linewidth]{"BL_25.jpg"}

\subsubsection{Guanyuanshu (BL 26): Quan Nguyên Du}
\includegraphics[width=\linewidth]{"BL_26.jpg"}

\subsubsection{Xiaochangshu (BL 27): Tiểu Trường Du}
\includegraphics[width=\linewidth]{"BL_27.jpg"}

\subsubsection{Pangguangshu (BL 28): Bàng Quang Du}
\includegraphics[width=\linewidth]{"BL_28.jpg"}

\subsubsection{Zhonglushu (BL 29): Trung Lữ Du}
\includegraphics[width=\linewidth]{"BL_29.jpg"}

\subsubsection{Baihuanshu (BL 30): Bạch Hoàn Du}
\includegraphics[width=\linewidth]{"BL_30.jpg"}

\subsubsection{Shangliao (BL 31): Thượng Liêu}
\includegraphics[width=\linewidth]{"BL_31.jpg"}

\subsubsection{Ciliao (BL 32): Thứ Liêu}
\includegraphics[width=\linewidth]{"BL_32.jpg"}

\subsubsection{Zhongliao (BL 33): Trung Liêu}
\includegraphics[width=\linewidth]{"BL_33.jpg"}

\subsubsection{Xialiao (BL 34): Hạ Liêu}
\includegraphics[width=\linewidth]{"BL_34.jpg"}

\subsubsection{Huiyang (BL 35): Hội Dương}
\includegraphics[width=\linewidth]{"BL_35.jpg"}

\subsubsection{Chengfu (BL 36): Thừa Phủ}
\includegraphics[width=\linewidth]{"BL_36.jpg"}

\subsubsection{Yinmen (BL 37): Ân Môn}
\includegraphics[width=\linewidth]{"BL_37.jpg"}

\subsubsection{Fuxi (BL 38): Phù Khích}
\includegraphics[width=\linewidth]{"BL_38.jpg"}

\subsubsection{Weiyang (BL 39): Uỷ Dương}
\includegraphics[width=\linewidth]{"BL_39.jpg"}

\subsubsection{Weizhong (BL 40): Uỷ Trung}
\includegraphics[width=\linewidth]{"BL_40.jpg"}

\subsubsection{Fufen (BL 41): Phụ Phân}
\includegraphics[width=\linewidth]{"BL_41.jpg"}

\subsubsection{Pohu (BL 42): Phách Hộ}
\includegraphics[width=\linewidth]{"BL_42.jpg"}

\subsubsection{Gaohuangshu (BL 43): Cao Hoang}
\includegraphics[width=\linewidth]{"BL_43.jpg"}

\subsubsection{Shentang (BL 44): Thần Đường}
\includegraphics[width=\linewidth]{"BL_44.jpg"}

\subsubsection{Yixi (BL 45): Y Hy}
\includegraphics[width=\linewidth]{"BL_45.jpg"}

\subsubsection{Geguan (BL 46): Cách Quan}
\includegraphics[width=\linewidth]{"BL_46.jpg"}

\subsubsection{Hunmen (BL 47): Hồn Môn}
\includegraphics[width=\linewidth]{"BL_47.jpg"}

\subsubsection{Yanggang (BL 48): Dương Cương}
\includegraphics[width=\linewidth]{"BL_48.jpg"}

\subsubsection{Yishe (BL 49): Ý Xá}
\includegraphics[width=\linewidth]{"BL_49.jpg"}

\subsubsection{Weicang (BL 50): Vị Thương}
\includegraphics[width=\linewidth]{"BL_50.jpg"}

\subsubsection{Huangmen (BL 51): Hoang Môn}
\includegraphics[width=\linewidth]{"BL_51.jpg"}

\subsubsection{Zhishi (BL 52): Chí Thất}
\includegraphics[width=\linewidth]{"BL_52.jpg"}

\subsubsection{Baohuang (BL 53): Bào Hoang}
\includegraphics[width=\linewidth]{"BL_53.jpg"}

\subsubsection{Zhibian (BL 54): Trật Biên}
\includegraphics[width=\linewidth]{"BL_54.jpg"}

\subsubsection{Heyang (BL 55): Hợp Dương}
\includegraphics[width=\linewidth]{"BL_55.jpg"}

\subsubsection{Chengjin (BL 56): Thừa Cân}
\includegraphics[width=\linewidth]{"BL_56.jpg"}

\subsubsection{Chengshan (BL 57): Thừa Sơn}
\includegraphics[width=\linewidth]{"BL_57.jpg"}

\subsubsection{Feiyang (BL 58): Phi Dương}
\includegraphics[width=\linewidth]{"BL_58.jpg"}

\subsubsection{Fuyang (BL 59): Phụ Dương}
\includegraphics[width=\linewidth]{"BL_59.jpg"}

\subsubsection{Kunlun (BL 60): Côn Lôn}
\includegraphics[width=\linewidth]{"BL_60.jpg"}

\subsubsection{Pucan (BL 61): Bộc Tham}
\includegraphics[width=\linewidth]{"BL_61.jpg"}

\subsubsection{Shenmai (BL 62): Thân Mạch}
\includegraphics[width=\linewidth]{"BL_62.jpg"}

\subsubsection{Jinmen (BL 63): Kim Môn}
\includegraphics[width=\linewidth]{"BL_63.jpg"}

\subsubsection{Jinggu (BL 64): Kinh Cốt}
\includegraphics[width=\linewidth]{"BL_64.jpg"}

\subsubsection{Shugu (BL 65): Thúc Cốt}
\includegraphics[width=\linewidth]{"BL_65.jpg"}

\subsubsection{Zutonggu (BL 66): Thông Cốc}
\includegraphics[width=\linewidth]{"BL_66.jpg"}

\subsubsection{Zhiyin (BL 67): Chí Âm}
\includegraphics[width=\linewidth]{"BL_67_2.jpg"}

% \clearpage
\subsection{Kidney Meridian (KI): Kinh túc thiếu âm Thận}

\subsubsection{Yongquan (KI 1): Dũng Tuyền}
\includegraphics[width=\linewidth]{"KI_1.jpg"}

\subsubsection{Rangu (KI 2): Nhiên Cốc}
\includegraphics[width=\linewidth]{"KI_2.jpg"}

\subsubsection{Taixi (KI 3): Thái Khê}
\includegraphics[width=\linewidth]{"KI_3.jpg"}

\subsubsection{Dazhong (KI 4): Đại Chung}
\includegraphics[width=\linewidth]{"KI_4.jpg"}

\subsubsection{Shuiquan (KI 5): Thủy Tuyền}
\includegraphics[width=\linewidth]{"KI_5.jpg"}

\subsubsection{Zhaohai (KI 6): Chiếu Hải}
\includegraphics[width=\linewidth]{"KI_6.jpg"}

\subsubsection{Fuliu (KI 7): Phục Lưu}
\includegraphics[width=\linewidth]{"KI_7.jpg"}

\subsubsection{Jiaoxin (KI 8): Giao Tín}
\includegraphics[width=\linewidth]{"KI_8.jpg"}

\subsubsection{Zhubin (KI 9): Trúc Tân}
\includegraphics[width=\linewidth]{"KI_9.jpg"}

\subsubsection{Yingu (KI 10): Âm Cốc}
\includegraphics[width=\linewidth]{"KI_10.jpg"}

\subsubsection{Henggu (KI 11): Hoành Cốt}
\includegraphics[width=\linewidth]{"KI 11.jpg"}

\subsubsection{Dahe (KI 12): Đại Hách}
\includegraphics[width=\linewidth]{"KI 12.jpg"}

\subsubsection{Qixue (KI 13): Khí Huyệt}
\includegraphics[width=\linewidth]{"KI 13.jpg"}

\subsubsection{Siman (KI 14): Tứ Mãn}
\includegraphics[width=\linewidth]{"KI 14.jpg"}

\subsubsection{Zhongzhu (KI 15): Trung Chú}
\includegraphics[width=\linewidth]{"KI 15.jpg"}

\subsubsection{Huangshu (KI 16): Hoang Du}
\includegraphics[width=\linewidth]{"KI 16.jpg"}

\subsubsection{Shangqu (KI 17): Thương Khúc}
\includegraphics[width=\linewidth]{"KI 17.jpg"}

\subsubsection{Shiguan (KI 18): Thạch Quan}
\includegraphics[width=\linewidth]{"KI 18.jpg"}

\subsubsection{Yindu (KI 19): Âm Đô}
\includegraphics[width=\linewidth]{"KI 19.jpg"}

\subsubsection{Futonggu (KI 20): Phúc Thông Cốc}
\includegraphics[width=\linewidth]{"KI 20.jpg"}

\subsubsection{Youmen (KI 21): U Môn}
\includegraphics[width=\linewidth]{"KI 21.jpg"}

\subsubsection{Bulang (KI 22): Bộ Lang}
\includegraphics[width=\linewidth]{"KI 22.jpg"}

\subsubsection{Shenfeng (KI 23): Thần Phong}
\includegraphics[width=\linewidth]{"KI 23.jpg"}

\subsubsection{Lingxu (KI 24): Linh Khu}
\includegraphics[width=\linewidth]{"KI 24.jpg"}

\subsubsection{Shencang (KI 25): Thần Tàng}
\includegraphics[width=\linewidth]{"KI 25.jpg"}

\subsubsection{Yuzhong (KI 26): Hoắc Trung}
\includegraphics[width=\linewidth]{"KI 26.jpg"}

\subsubsection{Shufu (KI 27): Du Phủ}
\includegraphics[width=\linewidth]{"KI 27.jpg"}

% \clearpage
\subsection{Pericardium Meridian(PC): Kinh thủ quyết âm Tâm bào}
\subsubsection{Tianchi (PC 1): Thiên Trì}
\includegraphics[width=\linewidth]{"PC_1.jpg"}

\subsubsection{Tianquan (PC 2): Thiên Tuyền}
\includegraphics[width=\linewidth]{"PC_2.jpg"}

\subsubsection{Quze (PC 3): Khúc Trạch}
\includegraphics[width=\linewidth]{"PC_3.jpg"}

\subsubsection{Ximen (PC 4): Khích Môn}
\includegraphics[width=\linewidth]{"PC_4.jpg"}

\subsubsection{Jianshi (PC 5): Giản Sử}
\includegraphics[width=\linewidth]{"PC_5.jpg"}

\subsubsection{Neiguan (PC 6): Nội Quan}
\includegraphics[width=\linewidth]{"PC_6.jpg"}

\subsubsection{Daling (PC 7): Đại Lăng}
\includegraphics[width=\linewidth]{"PC_7.jpg"}

\subsubsection{Laogong (PC 8): Lao Cung}
\includegraphics[width=\linewidth]{"PC_8.jpg"}

\subsubsection{Zhongchong (PC 9): Trung Xung}
\includegraphics[width=\linewidth]{"PC_9.jpg"}

% \clearpage
\subsection{Hand Shaoyang (TH): Kinh thủ thiếu dương Tam Tiêu}

\subsubsection{Guanchong (TH 1): Quan Xung}
\includegraphics[width=\linewidth]{"TH_1.jpg"}

\subsubsection{Yemen (TH 2): Dịch Môn}
\includegraphics[width=\linewidth]{"TH_2.jpg"}

\subsubsection{Zhongzhu (TH 3): Trung Chử}
\includegraphics[width=\linewidth]{"TH_3.jpg"}

\subsubsection{Yangchi (TH 4): Dương Trì}
\includegraphics[width=\linewidth]{"TH_4.jpg"}

\subsubsection{Waiguan (TH 5): Ngoại Quan}
\includegraphics[width=\linewidth]{"TH_5.jpg"}

\subsubsection{Zhigou (TH 6): Chi Câu}
\includegraphics[width=\linewidth]{"TH_6.jpg"}

\subsubsection{Huizong (TH 7): Hội Tông}
\includegraphics[width=\linewidth]{"TH_7.jpg"}

\subsubsection{Sanyangluo (TH 8): Tam Dương Lạc}
\includegraphics[width=\linewidth]{"TH_8.jpg"}

\subsubsection{Sidu (TH 9): Tứ Độc}
\includegraphics[width=\linewidth]{"TH_9.jpg"}

\subsubsection{Tianjing (TH 10): Thiên Tỉnh}
\includegraphics[width=\linewidth]{"TH_10.jpg"}

\subsubsection{Qinglengyuan (TH 11): Thanh Lãnh Uyên}
\includegraphics[width=\linewidth]{"TH_11.jpg"}

\subsubsection{Xiaoluo (TH 12): Tiêu Lạc}
\includegraphics[width=\linewidth]{"TH_12.jpg"}

\subsubsection{Naohui (TH 13): Nhu Hội}
\includegraphics[width=\linewidth]{"TH_13.jpg"}

\subsubsection{Jianliao (TH 14): Kiên Liêu}
\includegraphics[width=\linewidth]{"TH_14.jpg"}

\subsubsection{Tianliao (TH 15): Thiên Liêu}
\includegraphics[width=\linewidth]{"TH_15.jpg"}

\subsubsection{Tianyou (TH 16): Thiên Dũ}
\includegraphics[width=\linewidth]{"TH 16.jpg"}

\subsubsection{Yifeng (TH 17): Ế Phong}
\includegraphics[width=\linewidth]{"TH 17.jpg"}

\subsubsection{Chimai (TH 18): Khế Mạch}
\includegraphics[width=\linewidth]{"TH 18.jpg"}

\subsubsection{Luxi (TH 19): Lư Tức}
\includegraphics[width=\linewidth]{"TH 19.jpg"}

\subsubsection{Jiaosun (TH 20): Giác Tôn}
\includegraphics[width=\linewidth]{"TH 20.jpg"}

\subsubsection{Ermen (TH 21): Nhĩ Môn}
\includegraphics[width=\linewidth]{"TH 21.jpg"}

\subsubsection{Erheliao (TH 22): Nhĩ Hòa Liêu}
\includegraphics[width=\linewidth]{"TH 22.jpg"}

\subsubsection{Sizhukong (TH 23): Ty Trúc Không}
\includegraphics[width=\linewidth]{"TH 23.jpg"}
Nằm ở chỗ lõm tại đuôi lông mày (đầu ngoài của lông mày), ngay trên bờ trên hốc mắt.
% \clearpage
\subsection{Gallbladder Meridian (GB): Kinh túc thiếu dương Đởm}
\subsubsection{Tongziliao (GB 1): Đồng Tử Liêu}
\includegraphics[width=\linewidth]{"GB 1.jpg"}
\includegraphics[width=\linewidth]{"GB 1 (2).jpg"}
Huyệt này nằm cách khóe mắt ngoài 0,5 thốn về phía bên (ngoài),
 nằm trong chỗ lõm phía ngoài của bờ xương hốc mắt.


\subsubsection{Tinghui (GB 2): Thính Hội}
\includegraphics[width=\linewidth]{"GB 2.jpg"}

\subsubsection{Shangguan (GB 3): Thượng Quan}
\includegraphics[width=\linewidth]{"GB 3_0.jpg"}

\subsubsection{Hanyan (GB 4): Hàm Yến}
\includegraphics[width=\linewidth]{"GB 4.jpg"}

\subsubsection{Xuanlu (GB 5): Huyền Lư}
\includegraphics[width=\linewidth]{"GB 5.jpg"}

\subsubsection{Xuanli (GB 6): Huyền Ly}
\includegraphics[width=\linewidth]{"GB 6.jpg"}

\subsubsection{Qubin (GB 7): Khúc Tân}
\includegraphics[width=\linewidth]{"GB 7.jpg"}

\subsubsection{Shuaigu (GB 8): Suất Cốc}
\includegraphics[width=\linewidth]{"GB 8.jpg"}

\subsubsection{Tianchong (GB 9): Thiên Xung}
\includegraphics[width=\linewidth]{"GB 9.jpg"}

\subsubsection{Fubai (GB 10): Phù Bạch}
\includegraphics[width=\linewidth]{"GB 10.jpg"}

\subsubsection{Touqiaoyin (GB 11): Đầu Khiếu Âm}
\includegraphics[width=\linewidth]{"GB 11.jpg"}

\subsubsection{Wangu (GB 12): Hoàn Cốt}
\includegraphics[width=\linewidth]{"GB 12.jpg"}

\subsubsection{Benshen (GB 13): Bản Thần}
\includegraphics[width=\linewidth]{"GB 13.jpg"}

\subsubsection{Yangbai (GB 14): Dương Bạch}
\includegraphics[width=\linewidth]{"GB 14.jpg"}

\subsubsection{Toulinqi (GB 15): Đầu Lâm Khấp}
\includegraphics[width=\linewidth]{"GB 15.jpg"}
\includegraphics[width=\linewidth]{"GB 15 (2).jpg"}

\subsubsection{Muchuang (GB 16): Mục Song}
\includegraphics[width=\linewidth]{"GB 16.jpg"}

\subsubsection{Zhengying (GB 17): Chính Doanh}
\includegraphics[width=\linewidth]{"GB 17.jpg"}

\subsubsection{Chengling (GB 18): Thừa Linh}
\includegraphics[width=\linewidth]{"GB 18.jpg"}

\subsubsection{Naokong (GB 19): Não Không}
\includegraphics[width=\linewidth]{"GB19.jpg"}

\subsubsection{Fengchi (GB 20): Phong Trì}
\includegraphics[width=\linewidth]{"GB20.jpg"}

\subsubsection{Jianjing (GB 21): Kiên Tỉnh}
\includegraphics[width=\linewidth]{"GB 21.jpg"}

\subsubsection{Yuanye (GB 22): Uyển Dịch}
\includegraphics[width=\linewidth]{"GB22.jpg"}

\subsubsection{Zhejin (GB 23): Triệp Cân}
\includegraphics[width=\linewidth]{"GB23.jpg"}

\subsubsection{Riyue (GB 24): Nhật Nguyệt}
\includegraphics[width=\linewidth]{"GB 24.jpg"}

\subsubsection{Jingmen (GB 25): Kinh Môn}
\includegraphics[width=\linewidth]{"GB25.jpg"}

\subsubsection{Daimai (GB 26): Đới Mạch}
\includegraphics[width=\linewidth]{"GB 26.jpg"}

\subsubsection{Wushu (GB 27): Ngũ Khu}
\includegraphics[width=\linewidth]{"GB27.jpg"}

\subsubsection{Weidao (GB 28): Duy Đạo}
\includegraphics[width=\linewidth]{"GB28.jpg"}

\subsubsection{Juliao (GB 29): Cự Liêu}
\includegraphics[width=\linewidth]{"GB_29.jpg"}

\subsubsection{Huantiao (GB 30): Hoàn Khiêu}
\includegraphics[width=\linewidth]{"GB_30.jpg"}

\subsubsection{Fengshi (GB 31): Phong Thị}
\includegraphics[width=\linewidth]{"GB_31.jpg"}

\subsubsection{Zhongdu (GB 32): Trung Độc}
\includegraphics[width=\linewidth]{"GB_32.jpg"}

\subsubsection{Xiyangguan (GB 33): Tất Dương Quan}
\includegraphics[width=\linewidth]{"GB_33.jpg"}

\subsubsection{Yanglingquan (GB 34): Dương Lăng Tuyền}
\includegraphics[width=\linewidth]{"GB_34.jpg"}

\subsubsection{Yangjiao (GB 35): Dương Giao}
\includegraphics[width=\linewidth]{"GB_35.jpg"}

\subsubsection{Waiqiu (GB 36): Ngoại Khâu}
\includegraphics[width=\linewidth]{"GB_36.jpg"}

\subsubsection{Guangming (GB 37): Quang Minh}
\includegraphics[width=\linewidth]{"GB_37.jpg"}

\subsubsection{Yangfu (GB 38): Dương Phụ}
\includegraphics[width=\linewidth]{"GB_38.jpg"}

\subsubsection{Xuanzhong (GB 39): Huyền Chung}
\includegraphics[width=\linewidth]{"GB_39.jpg"}

\subsubsection{Qiuxu (GB 40): Khâu Khư}
\includegraphics[width=\linewidth]{"GB_40.jpg"}

\subsubsection{Zulinqi (GB 41): Túc Lâm Khấp}
\includegraphics[width=\linewidth]{"GB_41.jpg"}

\subsubsection{Diwuhui (GB 42): Địa Ngũ Hội}
\includegraphics[width=\linewidth]{"GB_42.jpg"}

\subsubsection{Xiaxi (GB 43): Hiệp Khê}
\includegraphics[width=\linewidth]{"GB_43.jpg"}

\subsubsection{Zuqiaoyin (GB 44): Túc Khiếu Âm}
\includegraphics[width=\linewidth]{"GB_44.jpg"}
% \clearpage
\subsection{Foot Jueyin (LR): Kinh túc quyết âm Can}

\subsubsection{Dadun (LR 1): Đại Đốn}
\includegraphics[width=\linewidth]{"LR_1.jpg"}

\subsubsection{Xingjian (LR 2): Hành Gian}
\includegraphics[width=\linewidth]{"LR_2.jpg"}

\subsubsection{Taichong (LR 3): Thái Xung}
\includegraphics[width=\linewidth]{"LR_3.jpg"}

\subsubsection{Zhongfeng (LR 4): Trung Phong}
\includegraphics[width=\linewidth]{"LR 4.jpg"}

\subsubsection{Ligou (LR 5): Lãi Câu}
\includegraphics[width=\linewidth]{"LR_5.jpg"}

\subsubsection{Zhongdu (LR 6): Trung Đô}
\includegraphics[width=\linewidth]{"LR_6.jpg"}

\subsubsection{Xiguan (LR 7): Tất Quan}
\includegraphics[width=\linewidth]{"LR_7.jpg"}

\subsubsection{Ququan (LR 8): Khúc Tuyền}
\includegraphics[width=\linewidth]{"LR_8.jpg"}

\subsubsection{Yinbao (LR 9): Âm Bao}
% \includegraphics[width=\linewidth]{"LR_9.jpg"}

\subsubsection{Zuwuli (LR 10): Túc Ngũ Lý}
% \includegraphics[width=\linewidth]{"LR_10.jpg"}

\subsubsection{Yinlian (LR 11): Âm Liêm}
% \includegraphics[width=\linewidth]{"LR_11.jpg"}

\subsubsection{Jimai (LR 12): Cấp Mạch}
\includegraphics[width=\linewidth]{"LR 12.jpg"}

\subsubsection{Zhangmen (LR 13): Chương Môn}
\includegraphics[width=\linewidth]{"LR 13.jpg"}

\subsubsection{Qimen (LR 14): Kỳ Môn}
\includegraphics[width=\linewidth]{"LR 14.jpg"}

% \clearpage
\subsection{Du Channel (GV / DU): Mạch Đốc}

\subsubsection{Changqiang (GV 1): Trường Cường}
\includegraphics[width=\linewidth]{"GV_1.jpg"}

\subsubsection{Yaoshu (GV 2): Yêu Du}
\includegraphics[width=\linewidth]{"GV_2.jpg"}

\subsubsection{Yaoyangguan (GV 3): Yêu Dương Quan}
\includegraphics[width=\linewidth]{"GV_3.jpg"}

\subsubsection{Mingmen (GV 4): Mệnh Môn}
\includegraphics[width=\linewidth]{"GV_4.jpg"}

\subsubsection{Xuanshu (GV 5): Huyền Khu}
\includegraphics[width=\linewidth]{"GV_5.jpg"}

\subsubsection{Jizhong (GV 6): Tích Trung}
\includegraphics[width=\linewidth]{"GV_6.jpg"}

\subsubsection{Zhongshu (GV 7): Trung Khu}
\includegraphics[width=\linewidth]{"GV_7.jpg"}

\subsubsection{Jinsuo (GV 8): Cân Súc}
\includegraphics[width=\linewidth]{"GV_8.jpg"}

\subsubsection{Zhiyang (GV 9): Chí Dương}
\includegraphics[width=\linewidth]{"GV_9.jpg"}

\subsubsection{Lingtai (GV 10): Linh Đài}
\includegraphics[width=\linewidth]{"GV_10.jpg"}

\subsubsection{Shendao (GV 11): Thần Đạo}
\includegraphics[width=\linewidth]{"GV_11_0.jpg"}

\subsubsection{Shenzhu (GV 12): Thân Trụ}
\includegraphics[width=\linewidth]{"GV_12.jpg"}

\subsubsection{Taodao (GV 13): Đào Đạo}
\includegraphics[width=\linewidth]{"GV_13.jpg"}

\subsubsection{Dazhui (GV 14): Đại Chùy}
\includegraphics[width=\linewidth]{"GV14.jpg"}
\includegraphics[width=\linewidth]{"GV_14.jpg"}

\subsubsection{Yamen (GV 15): Á Môn}
\includegraphics[width=\linewidth]{"GV15_0.jpg"}

\subsubsection{Fengfu (GV 16): Phong Phủ}
\includegraphics[width=\linewidth]{"GV16.jpg"}

\subsubsection{Naohu (GV 17): Na Hộ}
\includegraphics[width=\linewidth]{"GV17_0.jpg"}

\subsubsection{Qiangjian (GV 18): Cường Gian}
\includegraphics[width=\linewidth]{"GV18_0.jpg"}

\subsubsection{Houding (GV 19): Hậu Đỉnh}
\includegraphics[width=\linewidth]{"GV19.jpg"}

\subsubsection{Baihui (GV 20): Bách Hội}
\includegraphics[width=\linewidth]{"GV20_0.jpg"}
\includegraphics[width=\linewidth]{"GV20side.jpg"}

\subsubsection{Qianding (GV 21): Tiền Đỉnh}
\includegraphics[width=\linewidth]{"GV21.jpg"}

\subsubsection{Xinhui (GV 22): Tín Hội}
\includegraphics[width=\linewidth]{"GV22.jpg"}

\subsubsection{Shangxing (GV 23): Thượng Tinh}
\includegraphics[width=\linewidth]{"GV23.jpg"}

\subsubsection{Shenting (GV 24): Thần Đình}
\includegraphics[width=\linewidth]{"GV 24.jpg"}
\includegraphics[width=\linewidth]{"GV24.jpg"}

\subsubsection{Suliao (GV 25): Tố Liêu}
\includegraphics[width=\linewidth]{"GV 25.jpg"}

\subsubsection{Shuigou (GV 26): Thủy Câu (Nhân Trung)}
\includegraphics[width=\linewidth]{"GV 26.jpg"}

\subsubsection{Duiduan (GV 27): Đoài Đoan}
\includegraphics[width=\linewidth]{"GV 27.jpg"}

\subsubsection{Yinjiao (GV 28): Ngân Giao}
\includegraphics[width=\linewidth]{"GV 28.jpg"}

% \clearpage
\subsection{Conception Vessel (Ren Mai) (CV): Mạch Nhâm}
\subsubsection{Huiyin (CV 1): Hội Âm}
\includegraphics[width=\linewidth]{"CV 1.jpg"}

\subsubsection{Qugu (CV 2): Khúc Cốt}
\includegraphics[width=\linewidth]{"CV 2.jpg"}

\subsubsection{Zhongji (CV 3): Trung Cực}
\includegraphics[width=\linewidth]{"CV 3.jpg"}

\subsubsection{Guanyuan (CV 4): Quan Nguyên}
\includegraphics[width=\linewidth]{"CV 4.jpg"}

\subsubsection{Shimen (CV 5): Thạch Môn}
\includegraphics[width=\linewidth]{"CV 5.jpg"}

\subsubsection{Qihai (CV 6): Khí Hải}
\includegraphics[width=\linewidth]{"CV 6.jpg"}

\subsubsection{Yinjiao (CV 7): Âm Giao}
\includegraphics[width=\linewidth]{"CV 7.jpg"}

Ví trí: dưới rốn một tấc

Chỉ định: Viêm đường tiết niệu, viêm niêm mạc tử cung, đau sau sinh đẻ, đau bụng, mẩn ngứa.

\subsubsection{Shenque (CV 8): Thần Khuyết}
\includegraphics[width=\linewidth]{"CV 8.jpg"}

Vị trí: chính giữa rốn.

Chỉ định: Thường không châm, chỉ dùng để xác định.

\subsubsection{Shuifen (CV 9): Thủy Phân}
\includegraphics[width=\linewidth]{"CV 9.jpg"}

Vị trí: trên rốn một tấc

Chỉ định: đau bụng do viêm loét dạ dày - hành tá tràng, sưng mặt, bí đái.

\subsubsection{Xiawan (CV 10): Hạ Quản}
\includegraphics[width=\linewidth]{"CV 10.jpg"}

\subsubsection{Jianli (CV 11): Kiến Lý}
\includegraphics[width=\linewidth]{"CV 11.jpg"}

\subsubsection{Zhongwan (CV 12): Trung Quản}
\includegraphics[width=\linewidth]{"CV 12.jpg"}

Vị trí: giữa đường nối từ mũi kiếm xương ức đến rốn hoặc từ rốn thẳng lên 4 tấc.

Chỉ định: đau do viêm dạ dày - hành tá tràng, bụng trướng, tỳ vị bất hoà, ợ chua, ăn uống không tiêu.

\subsubsection{Shangwan (CV 13): Thượng Quản}
\includegraphics[width=\linewidth]{"CV 13.jpg"}

Vị trí: Từ rốn thẳng lên 5 tấc
Chỉ định: đau vùng ngực, nấc (vị khí nghịch), ợ chua, nôn mửa, nôn khan, điên cuồng.

\subsubsection{Juque (CV 14): Cự Khuyết}
\includegraphics[width=\linewidth]{"CV 14.jpg"}

\subsubsection{Jiuwei (CV 15): Cưu Vĩ}

\includegraphics[width=\linewidth]{"CV 15.jpg"}

\subsubsection{Zhongting (CV 16): Trung Đình}
\includegraphics[width=\linewidth]{"CV 16.jpg"}

\subsubsection{Shanzhong (CV 17): Đản Trung}
\includegraphics[width=\linewidth]{"CV 17.jpg"}

\subsubsection{Yutang (CV 18): Ngọc Đường}
\includegraphics[width=\linewidth]{"CV 18.jpg"}

\subsubsection{Zigong (CV 19): Tử Cung}
\includegraphics[width=\linewidth]{"CV 19.jpg"}

\subsubsection{Huagai (CV 20): Hoa Cái}
\includegraphics[width=\linewidth]{"CV 20.jpg"}

\subsubsection{Xuanji (CV 21): Toàn Cơ}
\includegraphics[width=\linewidth]{"CV 21.jpg"}

\subsubsection{Tiantu (CV 22): Thiên Đột}
\includegraphics[width=\linewidth]{"CV 22.jpg"}

\subsubsection{Lianquan (CV 23): Liêm Tuyền}
\includegraphics[width=\linewidth]{"CV 23.jpg"}

\subsubsection{Chengjiang (CV 24): Thừa Tương}
\includegraphics[width=\linewidth]{"CV 24.jpg"}
Nằm ở chỗ lõm dưới môi dưới và phía trên cằm, tại chính giữa rãnh môi - cằm.

% \subsection{Extra Acupuncture Points: Huyệt Ngoài Kinh (Kỳ Huyệt)}

\subsubsection{Sishencong (Ex-HN1): Tứ Thần Thông}
\includegraphics[width=\linewidth]{"HN1.jpg"}

\subsubsection{Dangyang (Ex-HN2): Đương Dương}
\includegraphics[width=\linewidth]{"Ex-HN2.jpg"}

\subsubsection{Yintang (Ex-HN3): Ấn Đường}
\includegraphics[width=\linewidth]{"Ex-HN3.jpg"}

\subsubsection{Yuyao (Ex-HN4): Ngư Yêu}
\includegraphics[width=\linewidth]{"Ex-HN4.jpg"}

\subsubsection{Taiyang (Ex-HN5): Thái Dương}
\includegraphics[width=\linewidth]{"Ex-HN5.jpg"}

\subsubsection{Erjian (Ex-HN6): Nhĩ Tiêm}
\includegraphics[width=\linewidth]{"Ex-HN6.jpg"}

\subsubsection{Qiuhou (Ex-HN7): Cầu Hậu}
\includegraphics[width=\linewidth]{"Ex-HN7.jpg"}

\subsubsection{Shangyingxiang (Ex-HN8): Thượng Nghinh Hương (Tỵ Thông)}
\includegraphics[width=\linewidth]{"Ex-HN8.jpg"}

\subsubsection{Neiyingxiang (Ex-HN9): Nội Nghinh Hương}
\includegraphics[width=\linewidth]{"Ex-HN9.jpg"}

\subsubsection{Juquan (Ex-HN10): Tụ Tuyền}
\includegraphics[width=\linewidth]{"HN10.jpg"}

\subsubsection{Haiquan (Ex-HN11): Hải Tuyền}
\includegraphics[width=\linewidth]{"HN11.jpg"}

\subsubsection{Jinjin (Ex-HN12): Kim Tân}
\includegraphics[width=\linewidth]{"HN12.jpg"}

\subsubsection{Yuye (Ex-HN13): Ngọc Dịch}
\includegraphics[width=\linewidth]{"HN13.jpg"}

\subsubsection{Yiming (Ex-HN14): Ế Minh}
\includegraphics[width=\linewidth]{"Ex_HN14.jpg"}

\subsubsection{Jingbailao (Ex-HN15): Cảnh Bách Lao}
\includegraphics[width=\linewidth]{"Ex-HN15.jpg"}

\subsubsection{Anmian (Ex-HN16): An Miên}
\includegraphics[width=\linewidth]{"HN16.jpg"}

\subsubsection{Bailao (Ex-HN17): Bách Lao}
\includegraphics[width=\linewidth]{"HN17.jpg"}

\subsubsection{Bitong (Ex-HN18): Tỵ Thông}
\includegraphics[width=\linewidth]{"Ex-HN18.jpg"}

\subsubsection{Jiachengjiang (Ex-HN19): Giáp Thừa Tương}
\includegraphics[width=\linewidth]{"Ex-HN19_0.jpg"}

\subsubsection{Qianzheng (Ex-HN20): Khiên Chính}
\includegraphics[width=\linewidth]{"Ex_HN20.jpg"}

\subsubsection{Shanglianquan (Ex-HN21): Thượng Liên Tuyền}
\includegraphics[width=\linewidth]{"Ex-HN21.jpg"}

\subsubsection{Jiali (Ex-HN22): Giáp Lý}
\includegraphics[width=\linewidth]{"HN22.jpg"}

\subsubsection{Xinshe (Ex-HN23): Tân Xá}
\includegraphics[width=\linewidth]{"HN23.jpg"}

\subsubsection{Xueyadian (Ex-HN24): Huyết Áp Điểm}
\includegraphics[width=\linewidth]{"HN24.jpg"}


% ==========================================
% 2. Chest and Abdomen Points (Vùng Ngực Bụng - CA)
% ==========================================
\subsubsection{Zigongxue (Ex-CA1): Tử Cung}
\includegraphics[width=\linewidth]{"Ex-CA1.jpg"}

\subsubsection{Qipang (Ex-CA2): Kỳ Bàng}
\includegraphics[width=\linewidth]{"Ex-CA2.jpg"}

\subsubsection{Tituo (Ex-CA3): Đề Thác}
\includegraphics[width=\linewidth]{"Ex-CA3.jpg"}

\subsubsection{Zhixie (Ex-CA4): Chỉ Tả}
\includegraphics[width=\linewidth]{"Ex-CA4.jpg"}

\subsubsection{Qizhongsibian (Ex-CA5): Kỳ Trung Tứ Biên}
\includegraphics[width=\linewidth]{"Ex-CA5.jpg"}

\subsubsection{Sanjiaojiu (Ex-CA6): Tam Giác Cứu}
\includegraphics[width=\linewidth]{"Ex-CA6.jpg"}

\subsubsection{Weishang (Ex-CA7): Vị Thượng}
\includegraphics[width=\linewidth]{"Ex-CA7.jpg"}


% ==========================================
% 3. Back Points (Vùng Lưng - B)
% ==========================================
\subsubsection{Dingchuan (Ex-B1): Định Suyễn (Suyễn Tức)}
\includegraphics[width=\linewidth]{"B1.jpg"}

\subsubsection{Huatuojiaji (Ex-B2): Hoa Đà Giáp Tích}
\includegraphics[width=\linewidth]{"Ex-B2.jpg"}

\subsubsection{Weiwanxiashu (Ex-B3): Vị Quản Hạ Du}
\includegraphics[width=\linewidth]{"Ex-B3.jpg"}

\subsubsection{Pigen (Ex-B4): Bĩ Căn}
\includegraphics[width=\linewidth]{"Ex-B4.jpg"}

\subsubsection{Xiazhishi (Ex-B5): Hạ Chí Thất}
\includegraphics[width=\linewidth]{"Ex-B5.jpg"}

\subsubsection{Yaoyi (Ex-B6): Yêu Nghi}
\includegraphics[width=\linewidth]{"Ex-B6.jpg"}

\subsubsection{Yaoyan (Ex-B7): Yêu Nhãn}
\includegraphics[width=\linewidth]{"Ex-B7.jpg"}

\subsubsection{Shiqizhuixia (Ex-B8): Thập Thất Chùy Hạ}
\includegraphics[width=\linewidth]{"Ex-B8.jpg"}

\subsubsection{Yaoqi (Ex-B9): Yêu Kỳ}
\includegraphics[width=\linewidth]{"Ex-B9.jpg"}

\subsubsection{Chuanxi (Ex-B10): Suyễn Tức}
\includegraphics[width=\linewidth]{"Ex-B10.jpg"}

\subsubsection{Huanmen (Ex-B11): Hoàn Môn}
\includegraphics[width=\linewidth]{"Ex-B11_0.jpg"}

\subsubsection{Huanzhong (Ex-B12): Hoàn Trung}
\includegraphics[width=\linewidth]{"Ex-B12.jpg"}

\subsubsection{Xiajiyu (Ex-B13): Hạ Tích Du}
\includegraphics[width=\linewidth]{"Ex-B13.jpg"}

\subsubsection{Juqueshu (Ex-B14): Cự Khuyết Du}
\includegraphics[width=\linewidth]{"Ex-B14_0.jpg"}

\subsubsection{Jieji (Ex-B15): Kết Tích}
\includegraphics[width=\linewidth]{"Ex-B15.jpg"}


% ==========================================
% 4. Upper Extremity Points (Vùng Chi Trên - UE)
% ==========================================
\subsubsection{Zhoujian (Ex-UE1): Trửu Tiêm}
\includegraphics[width=\linewidth]{"Ex-UE1.jpg"}

\subsubsection{Erbai (Ex-UE2): Nhị Bạch}
\includegraphics[width=\linewidth]{"Ex-UE2.jpg"}

\subsubsection{Zhongquan (Ex-UE3): Trung Tuyền}
\includegraphics[width=\linewidth]{"Ex-UE3.jpg"}

\subsubsection{Zhongkui (Ex-UE4): Trung Khôi}
\includegraphics[width=\linewidth]{"Ex-UE4.jpg"}

\subsubsection{Dagukong (Ex-UE5): Đại Cốt Không}
\includegraphics[width=\linewidth]{"Ex-UE5.jpg"}

\subsubsection{Xiaogukong (Ex-UE6): Tiểu Cốt Không}
\includegraphics[width=\linewidth]{"Ex-UE6.jpg"}

\subsubsection{Yaotongxue (Ex-UE7): Yêu Thống Điểm}
\includegraphics[width=\linewidth]{"Ex-UE7.jpg"}

\subsubsection{Wailaogong (Ex-UE8): Ngoại Lao Cung}
\includegraphics[width=\linewidth]{"Ex-UE8.jpg"}

\subsubsection{Baxie (Ex-UE9): Bát Tà}
\includegraphics[width=\linewidth]{"Ex-UE9.jpg"}

\subsubsection{Sifeng (Ex-UE10): Tứ Phùng}
\includegraphics[width=\linewidth]{"Ex-UE10.jpg"}

\subsubsection{Shixuan (Ex-UE11): Thập Tuyên}
\includegraphics[width=\linewidth]{"Ex-UE11.jpg"}

\subsubsection{Bizhong (Ex-UE12): Tý Trung}
\includegraphics[width=\linewidth]{"Ex-UE12.jpg"}

\subsubsection{Jianqian (Ex-UE13): Kiên Tiền}
\includegraphics[width=\linewidth]{"Ex-UE13.jpg"}

\subsubsection{Luozhen (Ex-UE14): Lạc Chẩm}
\includegraphics[width=\linewidth]{"Ex-UE14.jpg"}


% ==========================================
% 5. Lower Extremity Points (Vùng Chi Dưới - LE)
% ==========================================
\subsubsection{Kuangu (Ex-LE1): Khoan Cốt}
\includegraphics[width=\linewidth]{"Ex-LE1.jpg"}

\subsubsection{Heding (Ex-LE2): Hạc Đỉnh}
\includegraphics[width=\linewidth]{"Ex-LE2.jpg"}

\subsubsection{Xinei (Ex-LE3): Tất Nội}
\includegraphics[width=\linewidth]{"Ex-LE3.jpg"}

\subsubsection{Neixiyan (Ex-LE4): Nội Tất Nhãn}
\includegraphics[width=\linewidth]{"Ex-LE4.jpg"}

\subsubsection{Xiyan (Ex-LE5): Tất Nhãn}
\includegraphics[width=\linewidth]{"Ex-LE5.jpg"}

\subsubsection{Dannangxue (Ex-LE6): Đởm Nang}
\includegraphics[width=\linewidth]{"Ex-LE6_0.jpg"}

\subsubsection{Lanweixue (Ex-LE7): Lan Vĩ}
\includegraphics[width=\linewidth]{"Ex-LE7.jpg"}

\subsubsection{Neihuaijian (Ex-LE8): Nội Khỏa Tiêm}
\includegraphics[width=\linewidth]{"Ex-LE8_0.jpg"}

\subsubsection{Waihuaijian (Ex-LE9): Ngoại Khỏa Tiêm}
\includegraphics[width=\linewidth]{"Ex-LE9.jpg"}

\subsubsection{Bafeng (Ex-LE10): Bát Phong}
\includegraphics[width=\linewidth]{"Ex-LE10_0.jpg"}

\subsubsection{Duyin (Ex-LE11): Độc Âm}
\includegraphics[width=\linewidth]{"Ex-LE11.jpg"}

\subsubsection{Qiduan (Ex-LE12): Khí Đoan}
\includegraphics[width=\linewidth]{"Ex-LE12_0.jpg"}

\subsubsection{Baichongwo (Ex-LE13): Bách Trùng Oa}
\includegraphics[width=\linewidth]{"Ex-LE13.jpg"}

\subsubsection{Lineiting (Ex-LE14): Nội Lý Đình}
\includegraphics[width=\linewidth]{"Ex-LE14.jpg"}


% \subsection{EX-HN3 Ấn đường}
% \includegraphics[width=\linewidth]{"Ex-HN3.jpg"}

% Huyệt này nằm ở giữa hai đầu trong của lông mày, trên đường chính trung (đường dọc giữa).

% \subsection{GB14 Dương bạch}
% \includegraphics[width=\linewidth]{"GB 14.jpg"}

% Huyệt này nằm thẳng phía trên đồng tử (khi nhìn thẳng về phía trước) và cách điểm giữa lông mày 1 thốn về phía trên.

% \subsection{ST8 Đầu duy}
% \includegraphics[width=\linewidth]{"ST 8.jpg"}

% Huyệt này nằm phía trong đường chân tóc trước 0,5 thốn, cách DU-24 (Thần đình) 4,5 thốn về phía ngoài.

% Ghi chú: Huyệt này có thể được xác định tại giao điểm của một đường thẳng đứng giả định kẻ lùi về phía sau 0,5 thốn từ trong đường chân tóc thái dương và một đường nằm ngang kẻ vào 0,5 thốn từ trong đường chân tóc trước.

% \subsection{EX-HN5 Thái dương}
% \includegraphics[width=\linewidth]{"Ex-HN5.jpg"}

% Huyệt này nằm ở vùng thái dương, phía sau 1 thốn so với điểm giữa khóe mắt ngoài và đuôi lông mày.

% \subsection{GB4 Hàm yến}
% \includegraphics[width=\linewidth]{"GB 4.jpg"}

% Huyệt này nằm trong đường chân tóc tự nhiên, tại vị trí 1/4 khoảng cách giữa ST-8 (Đầu duy) và GB-7 (Khúc tân).

% \noindent
% \begin{center}
%     % Ảnh căn giữa dọc (align=c)
%     \includegraphics[align=c, width=0.45\textwidth]{"GB0701.png"}
%     \hfill
%     % Khối văn bản căn giữa dọc [c]
%     \begin{minipage}[c]{0.45\textwidth}
%         \subsection*{GB-7}
%         Vị trí huyệt: Huyệt này nằm trong đường chân tóc tự nhiên, cách SJ-20 (Giác tôn) một khoát ngón tay (trên cùng một đường nằm ngang).

%         Ghi chú: Huyệt này được xác định tốt nhất bằng cách ép/gập vành tai về phía trước (gập phần sau của gờ vành tai trên đè lên phần trước) rồi xác định huyệt SJ-20 (ngang mức với đỉnh tai), từ đó đo ngang ra một khoát ngón tay là huyệt GB-7 (Khúc tân).
%     \end{minipage}
% \end{center}

% \subsection{GV20 Bách hội}
% \includegraphics[width=\linewidth]{"GV20_0.jpg"}

% \noindent\includegraphics[width=\linewidth]{"GV20side.jpg"}


% Huyệt này nằm ở chỗ lõm, cách đường chân tóc trước 5 thốn và cách đường chân tóc sau 7 thốn.

% Ghi chú: Huyệt này cũng có thể được xác định vị trí là cách u chẩm ngoài 6 thốn và cách điểm giữa hai lông mày (glabella) 8 thốn.

% Ghi chú: Huyệt này nằm trên đường chính trung (đường dọc giữa).

% \subsection{EX-HN1 Tứ thần thông}

% \includegraphics[width=\linewidth]{"HN1.jpg"}

% Các huyệt này (tổng cộng 4 huyệt) nằm cách Bách hội DU-20 1 thốn về phía trước, phía sau và hai bên.

% \subsection{GV16 Phong phủ}
% \includegraphics[width=\linewidth]{"GV16.jpg"}

% Huyệt này nằm trong chỗ lõm ngay phía dưới u chẩm ngoài.

% Ghi chú: Huyệt này nằm cách mép đường chân tóc sau khoảng 1 thốn về phía trên.

% Ghi chú: U chẩm ngoài có thể được xác định bằng cách tìm điểm gờ xương nhô lên ở nền sọ trên đường chính trung (mặc dù điểm này có thể không cảm nhận rõ ở một số cá nhân).

% \begin{center}
%     \includegraphics[width=0.45\linewidth]{"DU1601.png"}
%     \hfill
%     \includegraphics[width=0.45\linewidth]{"DU1602.png"}
% \end{center}

% \subsection{BL10 Thiên trụ}
% \includegraphics[width=\linewidth]{"BL10.jpg"}

% Huyệt này nằm cách DU-15 (Á môn) 1,3 thốn về phía ngoài (ở bờ ngoài cơ thang).

% Ghi chú

% Xác định huyệt DU-16 (Phong phủ), huyệt nằm dưới u chẩm ngoài trong chỗ lõm được tạo bởi cơ thang hai bên.

% Sau đó xác định huyệt DU-15, huyệt nằm trong chỗ lõm dưới DU-16 1,5 thốn và cách đường chân tóc sau 0,5 thốn vào trong.

% Từ DU-15 đo ngang ra ngoài 1,3 thốn, bạn sẽ tìm thấy huyệt BL-10 (Thiên trụ).

% \begin{center}
%     \includegraphics[width=0.45\linewidth]{"BL1001.png"}
%     \hfill
%     \includegraphics[width=0.45\linewidth]{"BL1002.png"}
% \end{center}

% \subsection{GB20 Phong trì}
% \includegraphics[width=\linewidth]{"GB20.jpg"}

% Huyệt này nằm ở chỗ lõm giữa chỗ bám (nguyên ủy) của cơ ức đòn chũm và cơ thang.

% Ghi chú: Huyệt này nằm gần nền sọ (và thường nằm ở điểm đau/nhạy cảm nhất trong chỗ lõm).

% \begin{center}
%     \includegraphics[width=0.32\linewidth]{"CO_UC_DON_CHUM.png"}
%     \hfill
%     \includegraphics[width=0.32\linewidth]{"CO_THANG.png"}
%     \hfill
%     \includegraphics[width=0.32\linewidth]{"GB2001.png"}
% \end{center}

% % shoulder
% \subsection{SI17 Thiên dung}
% \includegraphics[width=\linewidth]{"SI 17.jpg"}

% Huyệt này nằm trong chỗ lõm ở giữa góc xương hàm dưới và bờ trước cơ ức đòn chũm.

% Ghi chú: Cơ ức đòn chũm dễ xác định hơn khi quay đầu sang một bên (đồng thời tác dụng lực cản ở cằm).

% Ghi chú: Huyệt này cũng có thể được xác định là nằm ở phía trước mỏm ngang của đốt sống C-2, nếu có thể sờ thấy.

% \begin{center}
%     \includegraphics[width=0.45\linewidth]{"SI1701.png"}
%     \hfill
%     \includegraphics[width=0.45\linewidth]{"SI1702.png"}
% \end{center}

% \subsection{EX-HN15 Cảnh bách lao}
% \includegraphics[width=\linewidth]{"Ex-HN15.jpg"}

% \subsection{Cảnh căn}

% Vị trí hình học: Bạn tìm giao điểm (góc lõm khoảng 45 độ) nơi đường bờ cổ đi thẳng xuống gặp đường ngang của bờ vai.

% Vị trí giải phẫu cơ: Huyệt nằm ngay trên khối cơ thang (cơ trapezius) bó trên, ngay vùng chân cổ sát với vai.

% \begin{enumerate}
%     \item Ngồi thẳng lưng, thả lỏng hoàn toàn hai vai.
%     \item Dùng ngón tay trỏ hoặc ngón tay giữa của tay đối diện (ví dụ: dùng tay phải tìm huyệt vai trái) sờ dọc theo đường bên cạnh cổ đi xuống dưới.
%     \item Khi ngón tay chạm đến điểm trũng lõm thấp nhất nơi chân cổ chuẩn bị bè ngang ra thành vai, đó chính là huyệt Cảnh Căn.
%     \item Khi ấn mạnh vào điểm này, bạn sẽ có cảm giác tê rát, tức nặng hoặc hơi buốt lan nhẹ lên cổ, chứng tỏ đã tìm đúng vị trí bó cơ căng cứng.
% \end{enumerate}

% \subsection{GB21 Kiên tỉnh}
% \includegraphics[width=\linewidth]{"GB 21.jpg"}

% Nằm ở điểm giữa nối từ huyệt Đại Chùy (Đại Chùy - GV 14) đến mỏm cùng vai, tại điểm cao nhất của vùng vai.

% \begin{center}
%     \includegraphics[width=0.45\linewidth]{"Acromion_of_left_scapula01.png"}
%     \hfill
%     \includegraphics[width=0.45\linewidth]{"Acromion_of_scapula06.png"}
% \end{center}

% \subsection{BL11 Đại trữ}

% \includegraphics[width=\linewidth]{"BL_11.jpg"}

% Huyệt này nằm cách bờ dưới mỏm gai của đốt sống ngực thứ nhất (T1) 1,5 thốn về phía bên.

% Ghi chú: Huyệt này nằm ở điểm cao nhất có thể nhìn thấy của khối cơ cạnh sống.

% \begin{center}
%     \includegraphics[width=0.45\linewidth]{"BL1101.png"}
%     \hfill
%     \includegraphics[width=0.45\linewidth]{"BL1102.png"}
% \end{center}


% \subsection{BL12 Phong môn}
% \includegraphics[width=\linewidth]{"BL_12.jpg"}

% \subsection{BL13 Phế du}
% \includegraphics[width=\linewidth]{"BL_13.jpg"}

% \subsection{BL14 Quyết âm du}
% \includegraphics[width=\linewidth]{"BL_14.jpg"}

% \subsection{BL43 Cao hoang}
% \includegraphics[width=\linewidth]{"BL_43.jpg"}

% \subsection{BL16 Đốc du}
% \includegraphics[width=\linewidth]{"BL_16.jpg"}

% \subsection{SI11 Thiên tông}
% \includegraphics[width=\linewidth]{"SI_11.jpg"}

% \subsection{SI9 Kiên trinh}
% \includegraphics[width=\linewidth]{"SI_9.jpg"}

% \subsection{SI10 Nhu du}
% \includegraphics[width=\linewidth]{"SI_10.jpg"}

% \subsection{LI15 Kiên ngung}
% \includegraphics[width=\linewidth]{"LI_15b.jpg"}
% % back

% \subsection{LI4 Hợp cốc}
% \includegraphics[width=\linewidth]{"LI_4.jpg"}

% Vị trí huyệt: Huyệt này nằm ở điểm chính giữa của xương bàn tay thứ hai, giữa xương bàn tay thứ nhất và xương bàn tay thứ hai.

% Ghi chú: Huyệt này có thể được xác định bằng cách khép ngón tay cái sát vào ngón tay trỏ, khi đó huyệt nằm ở vị trí cao nhất của khối cơ nổi lên, gần như ngang hàng với đầu nếp gấp (nếp hằn hổ khẩu).
% % Bấm chéo mắt cá đối diện

% \subsection{TH4 Dương trì}
% \includegraphics[width=\linewidth]{"TH_4.jpg"}
% % Chính giữa nếp lằn cổ tay phía sau (mặt ngoài)

% Vị trí huyệt: Huyệt này nằm ở chỗ lõm ngay giữa gân cơ duỗi ngón út và gân cơ duỗi chung các ngón, ngang mức khớp cổ tay.

% Ghi chú: Tìm khoảng kẽ giữa xương bàn tay thứ 4 và thứ 5, dùng ngón tay vuốt dọc theo khoảng kẽ này tới chỗ lõm ngang mức khớp cổ tay.

% Ghi chú: Các gân cơ sẽ nổi rõ hơn khi duỗi cổ tay và các ngón tay (khi thực hiện châm cứu, các gân cơ cần được thả lỏng).

% \subsection{TH5 Ngoại quan}
% \includegraphics[width=\linewidth]{"TH_5.jpg"}

% Vị trí huyệt: Huyệt này nằm ở chỗ lõm tạo thành giữa xương trụ và xương quay, ở phía xương quay của gân cơ duỗi chung các ngón tay, phía trên huyệt SJ-4 (Dương Trì) 2 thốn.

% Ghi chú: Huyệt nằm ở bờ xương quay, giữa xương quay và gân cơ duỗi chung các ngón.

% \subsection{LI11 Khúc trì}
% \includegraphics[width=\linewidth]{"LI_11.jpg"}

% Vị trí huyệt: Huyệt này nằm ở đầu ngoài của nếp gấp khuỷu tay, ở chính giữa khoảng cách giữa lồi cầu ngoài xương cánh tay và huyệt LU-5 (Xích Trạch).

% Ghi chú: Dễ xác định vị trí nhất khi gập khuỷu tay.

% \subsection{LI10 Thủ tam lý}
% % Khúc trì xuống 2 thốn (thường ghi nhận 2 thốn theo chuẩn Phân khu Học hội Trung Y; nếu theo tài liệu của bạn là 3 thốn)
% \includegraphics[width=\linewidth]{"LI_10.jpg"}

% Vị trí huyệt: Huyệt này nằm trên đường nối huyệt LI-5 (Dương Khê) đến huyệt LI-11 (Khúc Trì), phía trên huyệt LI-5 10 thốn và phía dưới huyệt LI-11 2 thốn.

% Ghi chú: Huyệt này nằm ở phía trên huyệt LI-9 (Thượng Liêm) 1 thốn.

% Ghi chú: Dễ xác định vị trí nhất khi gập khuỷu tay và hướng bờ xương quay của cẳng tay lên trên.

% \subsection{LI14 Tí nhu}
% \includegraphics[width=\linewidth]{"LI_14.jpg"}

% Vị trí huyệt: Huyệt này nằm ở chỗ lõm được tạo thành bởi đầu bám tận phía dưới của cơ delta và cơ cánh tay.

% Ghi chú: Khoảng chừng 3/5 khoảng cách tính từ huyệt LI-11 (Khúc Trì) đến huyệt LI-15 (Kiên Ngung).

% Ghi chú: Dễ xác định vị trí nhất khi gồng cứng các cơ ở cánh tay.

% \subsection{LI15 Kiên ngung}
% \includegraphics[width=\linewidth]{"LI_15b.jpg"}
% % Giơ ngang tay xuất hiện chỗ lõm
% Vị trí huyệt: Huyệt này nằm ở chỗ lõm tại nguyên ủy của cơ delta, phía trước và phía dưới mỏm cùng vai.

% Ghi chú: Dễ xác định vị trí nhất khi gồng cứng cánh tay, giúp làm lộ rõ nhất các điểm lõm (được tạo thành tại nguyên ủy của cơ delta).

% \subsection{GB21 Kiên tỉnh}
% \includegraphics[width=\linewidth]{"GB 21.jpg"}

% \subsection{SI9 Kiên trinh}
% \includegraphics[width=\linewidth]{"SI_9.jpg"}

% Vị trí huyệt: Khi thả lỏng cánh tay và ở tư thế khép sát vào thân người, huyệt này nằm ở chỗ lõm phía trên nếp gấp nách sau 1 thốn.
% Ghi chú: Huyệt này nằm ngay phía dưới bờ sau của cơ delta.

% \subsection{PC8 Lao cung}
% \includegraphics[width=\linewidth]{"PC_8.jpg"}

% Vị trí huyệt: Huyệt này nằm ở phía xương quay của xương bàn tay thứ 3, trong chỗ lõm giữa xương bàn tay thứ 2 và thứ 3.

% Ghi chú: Một cách khác để xác định huyệt này là nắm tay lại, huyệt nằm ngay tại vị trí mà đầu ngón tay giữa chạm vào lòng bàn tay.

% \subsection{LU10 Ngư tế}
% \includegraphics[width=\linewidth]{"LU_10.jpg"}

% \subsection{PC7 Đại lăng}
% \includegraphics[width=\linewidth]{"PC_7.jpg"}

% Vị trí huyệt: Huyệt này nằm ở giữa gân cơ gan tay dài và gân cơ gấp cổ tay quay, ngang mức huyệt HE-7 (Thần Môn).

% Ghi chú: Xác định huyệt HE-7 ở bờ gần xương đậu, hoặc xác định huyệt này nằm trên nếp gấp cổ tay.

% Ghi chú: Gân cơ gan tay dài có thể không xuất hiện ở một số người, trong trường hợp này hãy xác định huyệt nằm ở phía xương trụ của gân cơ gấp cổ tay quay.

% \subsection{PC6 Nội quan}
% \includegraphics[width=\linewidth]{"PC_6.jpg"}

% Vị trí huyệt: Huyệt này nằm ở giữa gân cơ gan tay dài và gân cơ gấp cổ tay quay, phía trên huyệt P-7 (Đại Lăng) 2 thốn.

% Ghi chú: Gân cơ gan tay dài có thể không xuất hiện ở một số người, trong trường hợp này hãy xác định huyệt nằm ở phía xương trụ của gân cơ gấp cổ tay quay.

% \subsection{LU5 Xích trạch}
% \includegraphics[width=\linewidth]{"LU_5.jpg"}

% Vị trí huyệt: Huyệt này nằm ở chỗ lõm phía xương quay của gân cơ nhị đầu cánh tay và nằm trên nếp gấp khuỷu tay.

% Ghi chú: Huyệt nằm hơi tách ra một chút, không nằm sát trực tiếp ngay cạnh gân cơ nhị đầu cánh tay.

% Ghi chú: Dễ xác định vị trí và thực hiện châm cứu nhất khi khuỷu tay hơi gập.

% \subsection{LU1 Trung phủ}
% \includegraphics[width=\linewidth]{"LU_1.jpg"}

% Vị trí huyệt: Huyệt này nằm ở khoang liên sườn thứ nhất, phía dưới huyệt LU-2 (Vân Môn) 1 thốn và cách đường giữa cơ thể 6 thốn.

% Ghi chú: Tam giác cơ delta - cơ ngực có thể dễ xác định hơn khi đưa tay về phía trước (khi áp dụng lực cản). Huyệt LU-2 nằm ở trung tâm của tam giác này, huyệt LU-1 nằm phía dưới 1 thốn và hơi chệch ra ngoài so với LU-2.

% Ghi chú: Trước tiên xác định sụn sườn của xương sườn thứ 2 (nằm ngang mức với góc xương ức), sau đó tìm điểm lõm của khoang liên sườn thứ nhất nằm ngay phía trên nó.


% \subsection{KI1 Dũng tuyền}
% \includegraphics[width=\linewidth]{"KI_1.jpg"}

% Vị trí huyệt: Huyệt này nằm ở lòng bàn chân, tại vị trí 1/3 khoảng cách tính từ đỉnh lòng bàn chân (gốc ngón chân thứ 2) đến bờ dưới gót chân, nằm giữa xương bàn chân thứ 2 và thứ 3.

% \begin{center}
%     \includegraphics[width=0.45\linewidth]{"KI0101.png"}
%     \hfill
%     \includegraphics[width=0.45\linewidth]{"KI0102.png"}
% \end{center}

% \subsection{BL60 Côn lôn}
% \includegraphics[width=\linewidth]{"BL_60.jpg"}

% Vị trí huyệt: Huyệt này nằm ở phía sau khớp cổ chân, trong chỗ lõm giữa đỉnh mắt cá ngoài và gân gót (gân Achilles).

% Ghi chú: Xác định ở trung tâm chỗ lõm, ở điểm chính giữa khoảng cách từ đỉnh mắt cá ngoài đến bờ sau của gân gót.

% \begin{center}
%     \includegraphics[width=0.45\linewidth]{"BL6001.png"}
%     \hfill
%     \includegraphics[width=0.45\linewidth]{"BL6002.png"}
% \end{center}

% \subsection{KI3 Thái khê}
% \includegraphics[width=\linewidth]{"KI_3.jpg"}

% Vị trí huyệt: Huyệt này nằm ở chỗ lõm ngang mức với đỉnh mắt cá trong và nằm giữa mắt cá trong với gân gót (gân Achilles).

% Ghi chú: Xác định huyệt ở chính giữa chỗ lõm nằm giữa đỉnh mắt cá trong và gân gót.

% \begin{center}
%     \includegraphics[width=0.45\linewidth]{"KI0301.png"}
%     \hfill
%     \includegraphics[width=0.45\linewidth]{"KI0302.png"}
% \end{center}

% \subsection{GB39 Huyền chung}
% \includegraphics[width=\linewidth]{"GB_39.jpg"}

% Vị trí huyệt: Huyệt này nằm ở giữa bờ sau xương mác và các gân cơ mác dài và mác ngắn, phía trên đỉnh mắt cá ngoài 3 thốn.

% Ghi chú: Huyệt này nằm ở vị trí phía trên đỉnh mắt cá ngoài một chiều rộng bàn tay (1 khoát tay).

% \begin{center}
%     \includegraphics[width=0.32\linewidth]{"GB3902.png"}
%     \hfill
%     \includegraphics[width=0.32\linewidth]{"GB3903.png"}
%     \hfill
%     \includegraphics[width=0.32\linewidth]{"GB3901.png"}
% \end{center}

% \subsection{SP6 Tam âm giao}
% \includegraphics[width=\linewidth]{"SP_6.jpg"}

% Vị trí huyệt: Huyệt này nằm ở trong chỗ lõm, phía trên đỉnh mắt cá trong 3 thốn, sát bờ trong xương chày.

% Ghi chú: Huyệt cũng có thể được xác định ở vị trí phía trên đỉnh mắt cá trong một chiều rộng bàn tay (3 thốn / 1 khoát tay).

% \begin{center}
%     \includegraphics[width=0.45\linewidth]{"SP0601.png"}
%     \hfill
%     \includegraphics[width=0.45\linewidth]{"SP0602.png"}
% \end{center}

% \subsection{BL57 Thừa sơn}
% \includegraphics[width=\linewidth]{"BL_57.jpg"}

% Huyệt này nằm ở trung điểm giữa BL-40 (Ủy Trung) và BL-60 (Côn Lôn), cách BL-40 8 thốn về phía dưới, tại chỗ lõm ngay dưới hai bụng cơ bụng chân (khi cơ co/gồng).

% Ghi chú

% Miết một ngón tay dọc theo gân Achilles (gân gót) cho đến khi cảm nhận được một chỗ lõm ngay dưới bụng cơ bụng chân (chỗ hợp lại của hai đầu cơ).

% Sẽ dễ tìm chỗ lõm này hơn khi cơ này gồng lên (ví dụ: bảo người bệnh đạp phần ức bàn chân vào tay bạn).

% \begin{center}
%     \includegraphics[width=0.45\linewidth]{"BL5701.png"}
%     \hfill
%     \includegraphics[width=0.45\linewidth]{"BL5702.png"}   
% \end{center}

% \subsection{BL40 Ủy trung}
% \includegraphics[width=\linewidth]{"BL_40.jpg"}

% Huyệt này nằm trên nếp gấp khoeo chân (nếp gấp nhượng chân), tại điểm lõm chính giữa gân cơ bán gân (semitendinosus) và gân cơ nhị đầu đùi (biceps femoris).

% Ghi chú

% Dễ tìm vị trí huyệt hơn khi khớp gối hơi gập (co nhẹ).

% \begin{center}
%     \includegraphics[width=0.32\linewidth]{"BL4001.png"}
%     \hfill
%     \includegraphics[width=0.32\linewidth]{"BL4002.png"}
%     \hfill
%     \includegraphics[width=0.32\linewidth]{"BL4003.png"}
% \end{center}

% \subsection{BL37 Ân môn}
% \includegraphics[width=\linewidth]{"BL_37.jpg"}

% Huyệt này nằm trên đường (tưởng tượng) nối giữa BL-40 (Ủy Trung) và BL-36 (Thừa Phù), nằm giữa các cơ gân đùi sau (hamstring muscles), dưới BL-36 6 thốn và trên BL-40 8 thốn.

% Ghi chú: Huyệt này có thể xác định bằng vị trí phía dưới BL-36 khoảng hai chiều rộng bàn tay (khoảng 6 thốn).

% Ghi chú: Huyệt này cũng có thể được xác định bằng cách chia đôi khoảng cách giữa BL-40 và BL-36, huyệt BL-37 sẽ nằm phía trên điểm trung điểm tham chiếu này 1 thốn.

% \begin{center}
%     \includegraphics[width=0.4\linewidth]{"BL3701.png"}
%     \hfill
%     \includegraphics[width=0.4\linewidth]{"BL3702.png"}
% \end{center}

% \subsection{BL36 Thừa phù}
% \includegraphics[width=\linewidth]{"BL_36.jpg"}

% Huyệt này nằm tại chỗ lõm giữa các cơ gân đùi sau (cơ gân kheo), ngay chính giữa nếp lằn mông (nếp gấp nếp mông), thẳng phía trên huyệt BL-40 (Ủy Trung).

% \begin{center}
%     \includegraphics[width=0.4\linewidth]{"BL3601.png"}
%     \hfill
%     \includegraphics[width=0.4\linewidth]{"BL3602.png"}
% \end{center}

% \subsection{BL54 Trật biên}
% % Gấp gót vuông góc, mũi gót ở đâu thì huyệt ở đó
% \includegraphics[width=\linewidth]{"BL_54.jpg"}

% Huyệt này nằm ở chỗ lõm đo ngang ra 3 thốn tính từ đường trung tuyến (đường giữa), ngang hàng với khe cùng-cụt (lỗ mỏm cùng-cụt).

% Ghi chú: Nếu thả lỏng vai, bờ trong của xương bả vai có thể được dùng làm mốc đo khoảng cách 3 thốn từ đường trung tuyến.

% Khe cùng-cụt nằm ở phía dưới mỏm gai đốt sống cùng thứ 4 (S4), nằm ở khoảng giữa sừng xương cùng và sừng xương cụt.

% \begin{center}
%     \includegraphics[width=0.4\linewidth]{"BL5401.png"}
%     \hfill
%     \includegraphics[width=0.4\linewidth]{"BL5402.png"}
% \end{center}

% \subsection{GB31 Phong thị}
% % Buông thõng tay tự do, đầu ngón tay giữa chạm ở đâu thì huyệt ở đó
% \includegraphics[width=\linewidth]{"GB_31.jpg"}

% Huyệt này nằm phía trên nếp gấp khoeo chân (nếp nhượng chân) 7 thốn và nằm thẳng phía dưới mấu chuyển lớn xương đùi.

% Ghi chú

% Có thể tìm mấu chuyển lớn bằng cách đặt một bàn tay lên khu vực nghi ngờ, đồng thời dùng tay kia xoay bàn chân để cảm nhận sự dịch chuyển của mấu chuyển lớn, giúp xác định vị trí chính xác hơn.

% Khoảng cách từ điểm lồi nhất của mấu chuyển lớn đến nếp gấp khoeo chân là 19 thốn; lấy một điểm tham chiếu tưởng tượng ở vị trí 2/3 khoảng cách này (đo từ mấu chuyển lớn xuống), huyệt GB-31 sẽ nằm phía trên điểm tham chiếu này khoảng 1 thốn.

% Ghi chú

% Huyệt này cũng được cho là nằm ngay đầu ngón tay giữa khi đứng thẳng, hai tay buông thõng áp sát dọc hai bên thân người (xuôi theo đùi); tuy nhiên, do tỷ lệ cơ thể ở mỗi người là khác nhau nên cách này không phải lúc nào cũng hoàn toàn chính xác.

% \begin{center}
%     \includegraphics[width=0.4\linewidth]{"GB3101.png"}
%     \hfill
%     \includegraphics[width=0.4\linewidth]{"GB3102.png"}    
% \end{center}

% \subsection{GB34 Dương lăng tuyền}
% \includegraphics[width=\linewidth]{"GB_34.jpg"}

% Huyệt này nằm tại chỗ lõm phía trước và phía dưới 1 thốn so với chỏm (đầu) xương mác.

% Ghi chú

% Miết ngón tay từ mặt ngoài cẳng chân hướng lên trên cho đến khi phần mềm của cơ được thay thế bằng gờ xương nhô lên của chỏm xương mác.

% \begin{center}
%     \includegraphics[width=0.4\linewidth]{"GB3401.png"}
%     \hfill
%     \includegraphics[width=0.4\linewidth]{"GB3403.png"}
% \end{center}

% \subsection{ST36 Túc tam lý}
% \includegraphics[width=\linewidth]{"ST_36.jpg"}

% Huyệt này nằm phía dưới huyệt ST-35 (Độc Tỵ) 3 thốn và đo ngang ra 1 khoát ngón tay (ngang 1 ngón tay) tính từ mào trước xương chày.

% Ghi chú: Huyệt này cũng có thể được xác định ở vị trí phía dưới huyệt GB-34 (Dương Lăng Tuyền) 1 thốn và dưới ST-35 khoảng 1 chiều rộng bàn tay (4 ngón tay khép lại / khoảng 3 thốn) (đồng thời cách mào trước xương chày 1 khoát ngón tay về phía ngoài).


% \begin{center}
%     \includegraphics[width=0.4\linewidth]{"ST3601.png"}
%     \hfill
%     \includegraphics[width=0.4\linewidth]{"ST3602.png"}
% \end{center}

% \subsection{EX-LE10 Bát phong}
% \includegraphics[width=\linewidth]{"Ex-LE10_0.jpg"}

% Vị trí huyệt: Các huyệt này (4 huyệt ở mỗi bên bàn chân) nằm ở giữa kẽ các ngón chân, cách mép kẽ ngón chân 0,5 thốn về phía cổ chân.

% Ghi chú: Tổng cộng gồm 8 huyệt, bao gồm cả các huyệt Hành Gian (LIV-2), Nội Đình (ST-44) và Hiệp Khê (GB-43).

% \begin{center}
%     \includegraphics[width=0.4\linewidth]{"EXLE1001.png"}
%     \hfill
%     \includegraphics[width=0.4\linewidth]{"EXLE1002.png"}
% \end{center}

% \subsection{LR3 Thái xung}
% \includegraphics[width=\linewidth]{"LR_3.jpg"}

% Huyệt này nằm ở phía trước (phía xa) so với giao điểm của xương bàn chân I và II (xương đốt bàn chân 1 và 2).

% Ghi chú: Để xác định huyệt này, miết ngón tay vào kẽ giữa ngón chân thứ nhất và ngón chân thứ hai bắt đầu từ huyệt LIV-2 (Hành Gian) hướng về phía cổ chân, cho đến khi thấy một chỗ lõm nằm ngay trước giao điểm của xương bàn chân I và II.

% \begin{center}
%     \includegraphics[width=0.4\linewidth]{"LR0301.png"}
%     \hfill
%     \includegraphics[width=0.4\linewidth]{"LR0302.png"}
% \end{center}

% \subsection{ST44 Nội đình}
% \includegraphics[width=\linewidth]{"ST_44.jpg"}

% Huyệt này nằm ở giữa xương bàn chân II và III (đốt bàn 2 và 3), phía trên (phía gần cổ chân) cách bờ kẽ ngón chân 0,5 thốn.

% \begin{center}
%     \includegraphics[width=0.4\linewidth]{"ST4401.png"}
%     \hfill
%     \includegraphics[width=0.4\linewidth]{"ST4402.png"}
% \end{center}

% \subsection{ST41 Giải khê}
% \includegraphics[width=\linewidth]{"ST 41.jpg"}

% Vị trí huyệt: Huyệt này nằm ở chỗ lõm giữa gân cơ duỗi ngón cái dài và gân cơ duỗi các ngón chân dài, nằm ngang mức với đỉnh mắt cá ngoài.

% Ghi chú: Gân cơ duỗi ngón cái dài sẽ nổi rõ hơn khi duỗi ngón chân cái (khi có lực cản), gân cơ duỗi các ngón chân sẽ nổi rõ hơn khi duỗi bốn ngón chân còn lại chống lại lực cản.

% \begin{center}
%     \includegraphics[width=0.4\linewidth]{"ST4101.png"}
%     \hfill
%     \includegraphics[width=0.4\linewidth]{"ST4102.png"}
% \end{center}

% \subsection{ST35 Độc tỵ}
% \includegraphics[width=\linewidth]{"ST_35.jpg"}

% Vị trí huyệt: Huyệt này nằm ở giữa bờ sau xương mác và các gân cơ mác dài và mác ngắn, phía trên đỉnh mắt cá ngoài 3 thốn.

% Ghi chú: Huyệt này nằm ở vị trí phía trên đỉnh mắt cá ngoài một chiều rộng bàn tay (1 khoát tay).

% \begin{center}
%     \includegraphics[width=0.32\linewidth]{"ST3503.png"}
%     \hfill
%     \includegraphics[width=0.32\linewidth]{"ST3501.png"}
%     \hfill
%     \includegraphics[width=0.32\linewidth]{"ST3502.png"}
% \end{center}

% \subsection{EX-LE5 Tất nhãn (Xiyan)}
% \includegraphics[width=\linewidth]{"Ex-LE5.jpg"}

% Vị trí huyệt: Các huyệt này (gồm 2 huyệt) nằm phía dưới xương bánh chè, trong các chỗ lõm được tạo thành khi gập đầu gối, ở hai bên dây chằng xương bánh chè.

% Ghi chú: Huyệt MN-LE-16 nằm ở cùng vị trí với huyệt Độc Tỵ (ST-35).

% \subsection{EX-LE2 Hạc đỉnh}
% \includegraphics[width=\linewidth]{"Ex-LE2.jpg"}

% Vị trí huyệt: Huyệt này nằm ở chỗ lõm tại điểm giữa bờ trên xương bánh chè.

% \begin{center}
%     \includegraphics[width=0.45\linewidth]{"EXLE201.png"}
%     \hfill
%     \includegraphics[width=0.45\linewidth]{"EXLE202.png"}
% \end{center}


% \subsection{SP10 Huyết hải}
% \includegraphics[width=\linewidth]{"SP_10a.jpg"}

% Vị trí huyệt: Huyệt này nằm ở chỗ lõm trên phần gồ của cơ rộng trong, nằm thẳng phía trên huyệt SP-9 (Âm Lăng Tuyền) và phía trên bờ trên xương bánh chè 2 thốn.

% Ghi chú: Huyệt này có thể được xác định bằng cách đặt lòng bàn tay lên bờ dưới xương bánh chè (ngón tay cái hướng vào trong), duỗi thẳng tất cả các ngón tay và mở ngón cái ra một góc 45° so với ngón trỏ, huyệt khi đó nằm ngay dưới đầu ngón tay cái.

% Ghi chú: Lấy khoảng cách bằng hai lần chiều cao xương bánh chè (chiều cao xương bánh chè là 2 thốn) tính hướng lên từ bờ trên xương bánh chè và nằm thẳng phía trên SP-9.

% \begin{center}
%     \includegraphics[width=0.32\linewidth]{"SP1001.png"}
%     \hfill
%     \includegraphics[width=0.32\linewidth]{"SP1002.png"}
%     \hfill
%     \includegraphics[width=0.32\linewidth]{"SP1003.png"}
% \end{center}

% \subsection{ST34 Lương khâu}
% \includegraphics[width=\linewidth]{"ST_34.jpg"}

% Huyệt này nằm tại chỗ lõm phía trên bờ trên xương bánh chè 2 thốn. Huyệt nằm trên đường (tưởng tượng) nối từ bờ ngoài xương bánh chè đến gai chậu trước trên.

% Ghi chú: Gấp đôi chiều cao của xương bánh chè (= 2 thốn) đo nối tiếp lên phía trên từ bờ trên của chính nó.

% \begin{center}
%     \includegraphics[width=0.45\linewidth]{"ST3401.png"}
%     \hfill
%     \includegraphics[width=0.45\linewidth]{"ST3402.png"}
% \end{center}

% % \subsection{GB14 Yang Bai}
% % Huyệt này nằm ngay phía trên đồng tử (khi nhìn thẳng về phía trước) và phía trên trung điểm của lông mày 1 thốn.

% % \begin{center}
% %     \includegraphics[width=0.45\linewidth]{"GB14YangBaiGlobal.png"}
% %     \hfill
% %     \includegraphics[width=0.45\linewidth]{"GB14YangBaiLocal.png"}
% % \end{center}

% \subsection{CV22 Thiên đột}
% \includegraphics[width=\linewidth]{"CV 22.jpg"}

% Vị trí huyệt: Nằm trên đường giữa cơ thể, ở chính giữa hố trên xương ức, phía trên khuyết trên xương ức 0,5 thốn (khoảng 1 cm).

% \begin{center}
%     \includegraphics[width=0.3\linewidth]{"Manabrium_sterni.png"}
%     \hfill
%     \includegraphics[width=0.3\linewidth]{"REN2201.png"}
%     \hfill
%     \includegraphics[width=0.3\linewidth]{"REN2202.png"}
% \end{center}

% \subsection{CV19 Tử cung}
% \includegraphics[width=1\linewidth]{"CV 19.jpg"}

% Vị trí huyệt: Nằm ngang mức với chỗ tiếp giáp của khoang liên sườn 2 và xương ức, trên đường giữa cơ thể.

% Ghi chú: Trước tiên xác định sụn sườn của xương sườn thứ 2, sau đó tìm điểm lõm của khoang liên sườn 2 nằm ngay phía dưới.

% \begin{center}
%     \includegraphics[width=0.3\linewidth]{"costa02.png"}
%     \hfill
%     \includegraphics[width=0.3\linewidth]{"REN1901.png"}
%     \hfill
%     \includegraphics[width=0.3\linewidth]{"REN1902.png"}
% \end{center}

% \subsection{CV17 Đản trung}
% \includegraphics[width=\linewidth]{"CV 17.jpg"}

% Vị trí huyệt: Nằm ngang mức với chỗ tiếp giáp của khoang liên sườn 4 và xương ức, trên đường giữa cơ thể.

% Ghi chú: Trước tiên xác định sụn sườn của xương sườn thứ 2, sau đó tìm điểm lõm của khoang liên sườn 2 nằm ngay phía dưới và đếm dần xuống.

% Ở nam giới, huyệt này có thể được xác định ngang mức với hai đầu núm vú, trên đường giữa cơ thể.


% \begin{center}
%     \includegraphics[width=0.2\linewidth]{"costa02.png"}
%     \hfill
%     \includegraphics[width=0.2\linewidth]{"costa04.png"}
%     \hfill
%     \includegraphics[width=0.2\linewidth]{"REN1703.png"}
%     \hfill
%     \includegraphics[width=0.2\linewidth]{"REN1702.png"}
% \end{center}

% \subsection{CV15 Cưu vĩ}
% \includegraphics[width=\linewidth]{"CV 15.jpg"}

% Vị trí huyệt: Huyệt này nằm ở phía trên rốn 7 thốn và phía dưới góc ức-sườn 1 thốn, trên đường giữa cơ thể.

% Ghi chú: Góc ức-sườn nằm ngang mức với phần xương ức tiếp giáp với mỏm mũi kiếm.

% Có thể tìm góc ức-sườn bằng cách vuốt ngón trỏ hướng về phía xương ức, dọc theo đường giữa, vào điểm lõm ngay bên dưới xương ức.

% Mỏm mũi kiếm nằm ở vị trí này; kích thước mỏm mũi kiếm thay đổi tùy theo từng người nên không phải lúc nào cũng có thể sờ hoặc nhìn thấy được.

% Ở người lớn tuổi, mỏm mũi kiếm có thể bị vôi hóa, khiến việc phân biệt xương ức với mỏm mũi kiếm trở nên khó khăn hơn.
% (Mỏm mũi kiếm có kích thước thay đổi tùy từng người, do đó huyệt này có thể nằm ngay trên mỏm mũi kiếm hoặc ở phía dưới mỏm mũi kiếm).


% \begin{center}
%     \includegraphics[width=0.2\linewidth]{"corpus_sterni.png"}
%     \hfill
%     \includegraphics[width=0.2\linewidth]{"processus_xiphoideus.png"}
%     \hfill
%     \includegraphics[width=0.2\linewidth]{"REN1501.png"}
%     \hfill
%     \includegraphics[width=0.2\linewidth]{"REN1502.png"}
% \end{center}

% \subsection{CV13 Thượng quản}
% \includegraphics[width=\linewidth]{"CV 13.jpg"}

% Vị trí huyệt: Huyệt này nằm ở phía trên rốn 5 thốn và phía dưới góc ức-sườn 3 thốn, trên đường giữa cơ thể.

% Ghi chú: Góc ức-sườn nằm ngang mức với phần xương ức tiếp giáp với mỏm mũi kiếm.

% Có thể tìm góc ức-sườn bằng cách vuốt ngón trỏ hướng về phía xương ức, dọc theo đường giữa, vào điểm lõm ngay bên dưới xương ức.

% Mỏm mũi kiếm nằm ở vị trí này; kích thước mỏm mũi kiếm thay đổi tùy theo từng người nên không phải lúc nào cũng có thể sờ hoặc nhìn thấy được.

% Ở người lớn tuổi, mỏm mũi kiếm có thể bị vôi hóa, khiến việc phân biệt xương ức với mỏm mũi kiếm trở nên khó khăn hơn.

% \begin{center}
%     \includegraphics[width=0.3\linewidth]{"corpus_sterni.png"}
%     \hfill
%     \includegraphics[width=0.3\linewidth]{"processus_xiphoideus.png"}
%     \hfill
%     \includegraphics[width=0.3\linewidth]{"REN1301.png"}
% \end{center}

% \subsection{CV12 Trung quản}
% \includegraphics[width=\linewidth]{"CV 12.jpg"}


% Vị trí huyệt: Huyệt này nằm ở phía trên rốn 4 thốn và ở chính giữa (1/2) khoảng cách giữa rốn và góc ức-sườn, trên đường giữa cơ thể.

% Ghi chú: Góc ức-sườn nằm ngang mức với phần xương ức tiếp giáp với mỏm mũi kiếm.

% Có thể tìm góc ức-sườn bằng cách vuốt ngón trỏ hướng về phía xương ức, dọc theo đường giữa, vào điểm lõm ngay bên dưới xương ức.

% Mỏm mũi kiếm nằm ở vị trí này; kích thước mỏm mũi kiếm thay đổi tùy theo từng người nên không phải lúc nào cũng có thể sờ hoặc nhìn thấy được.

% Ở người lớn tuổi, mỏm mũi kiếm có thể bị vôi hóa, khiến việc phân biệt xương ức với mỏm mũi kiếm trở nên khó khăn hơn.

% \begin{center}
%     \includegraphics[width=0.3\linewidth]{"corpus_sterni.png"}
%     \hfill
%     \includegraphics[width=0.3\linewidth]{"processus_xiphoideus.png"}
%     \hfill
%     \includegraphics[width=0.3\linewidth]{"REN1201.png"}
% \end{center}

% \subsection{CV10 Hạ quản}
% \includegraphics[width=\linewidth]{"CV 10.jpg"}

% Vị trí huyệt: Huyệt này nằm trên đường giữa cơ thể, phía trên rốn 2 thốn và phía dưới góc ức-sườn 6 thốn.

% Ghi chú: Góc ức-sườn nằm ngang mức với phần xương ức tiếp giáp với mỏm mũi kiếm.
% Có thể tìm góc ức-sườn bằng cách vuốt ngón trỏ hướng về phía xương ức, dọc theo đường giữa, vào điểm lõm ngay bên dưới xương ức.
% Mỏm mũi kiếm nằm ở vị trí này; kích thước mỏm mũi kiếm thay đổi tùy theo từng người nên không phải lúc nào cũng có thể sờ hoặc nhìn thấy được.
% Ở người lớn tuổi, mỏm mũi kiếm có thể bị vôi hóa, khiến việc phân biệt xương ức với mỏm mũi kiếm trở nên khó khăn hơn.


% \begin{center}
%     \includegraphics[width=0.3\linewidth]{"corpus_sterni.png"}
%     \hfill
%     \includegraphics[width=0.3\linewidth]{"processus_xiphoideus.png"}
%     \hfill
%     \includegraphics[width=0.3\linewidth]{"REN1001.png"}
% \end{center}


% \subsection{CV6 Khí hải}
% \includegraphics[width=\linewidth]{"CV 6.jpg"}

% Vị trí huyệt: Huyệt này nằm ở phía dưới rốn 1,5 thốn, phía trên bờ trên của khớp mu 3,5 thốn, trên đường giữa cơ thể.

% \begin{center}
%     \includegraphics[width=0.48\linewidth]{"REN0602.png"}
%     \hfill
%     \includegraphics[width=0.48\linewidth]{"REN0601.png"}
% \end{center}

% \subsection{CV3 Trung cực}
% \includegraphics[width=\linewidth]{"CV 3.jpg"}

% Vị trí huyệt: Huyệt này nằm ở phía dưới rốn 4 thốn, phía trên bờ trên của khớp mu 1 thốn, trên đường giữa cơ thể.

% \begin{center}
%     \includegraphics[width=0.48\linewidth]{"REN0301.png"}
%     \hfill
%     \includegraphics[width=0.48\linewidth]{"REN0302.png"}
% \end{center}


% \subsection{CV1 Hội âm}
% \includegraphics[width=\linewidth]{"CV 1.jpg"}

% Vị trí huyệt: Huyệt này nằm ở điểm chính giữa hậu môn và bìu ở nam giới, và giữa hậu môn và mép sau âm hộ ở nữ giới. Nằm trên đường giữa cơ thể.

% \begin{center}
%     \includegraphics[width=0.48\linewidth]{"REN0102.png"}
%     \hfill
%     \includegraphics[width=0.48\linewidth]{"REN0101.png"}
% \end{center}

% \subsection{GB26 Đới mạch}
% \includegraphics[width=\linewidth]{"GB 26.jpg"}

% Vị trí huyệt: Huyệt này nằm ngay phía dưới huyệt LIV-13 (Chương Môn - nằm ở phía trước và phía dưới đầu tự do của xương sườn thứ 11), ngang mức với rốn.

% Ghi chú: Để xác định đầu tự do của xương sườn thứ 12 (ở vùng thắt lưng) và xương sườn thứ 11 (thường nằm phía trên mức rốn), dùng tay vuốt dọc theo bờ dưới lồng ngực bắt đầu từ vùng bụng trên.
% Đầu tự do của xương sườn thứ 11 thường nằm trên đường nách giữa (đường thẳng đứng kẻ từ đỉnh nách).

% \begin{center}
%     \includegraphics[width=0.3\linewidth]{"GB2602.png"}
%     \hfill
%     \includegraphics[width=0.3\linewidth]{"GB2603.png"}
%     \hfill
%     \includegraphics[width=0.3\linewidth]{"GB2601.png"}
% \end{center}

% \subsection{SP15 Đại hoành}

% \includegraphics[width=\linewidth]{"SP 15.jpg"}

% Vị trí huyệt: Huyệt này nằm ngang mức với rốn.

% Ghi chú: Bờ ngoài của cơ thẳng bụng (có thể sờ thấy được) nằm cách đường giữa cơ thể 4 thốn.

% \begin{center}
%     \includegraphics[width=0.3\linewidth]{"SP1501.png"}
%     \hfill
%     \includegraphics[width=0.3\linewidth]{"SP1504.png"}
%     \hfill
%     \includegraphics[width=0.3\linewidth]{"SP1502.png"}
% \end{center}

% \subsection{Ex-CA1 Tử cung}
% \includegraphics[width=\linewidth]{"Ex-CA1.jpg"}

% Vị trí huyệt: Huyệt này nằm ở vùng bụng dưới, đo ngang ra 3 thốn từ huyệt Trung Cực (REN-3).


% \begin{center}
%     \includegraphics[width=0.45\linewidth]{"EXCA102.png"}
%     \hfill
%     \includegraphics[width=0.45\linewidth]{"EXCA101.png"}
% \end{center}

% \subsection{BL26 Quan nguyên du}
% \includegraphics[width=\linewidth]{"BL_26.jpg"}

% Vị trí huyệt: Huyệt này nằm đo ngang ra 1,5 thốn từ bờ dưới mỏm gai của đốt sống thắt lưng thứ 5 (L5).

% Ghi chú: Huyệt này nằm trên điểm nổi cao nhất có thể nhìn thấy của các cơ cạnh sống.

% \begin{center}
%     \includegraphics[width=0.3\linewidth]{"BL2601.png"}
%     \hfill
%     \includegraphics[width=0.3\linewidth]{"BL2603.png"}    
%     \hfill
%     \includegraphics[width=0.3\linewidth]{"BL2604.png"}        
% \end{center}

% % \section{Danh sách tất cả các huyệt đạo}
% % \includepdf[pages=-]{pdfs/danhsachhuyetdao.pdf}
